\documentclass[11pt,a4paper]{article}

\usepackage[margin=1in]{geometry}
\usepackage[T1]{fontenc}

\usepackage{lmodern}
\usepackage{amsmath,amssymb,bm}
\usepackage{booktabs}
\usepackage{longtable}
\usepackage{multirow}
\usepackage{array}
\usepackage{tabularx}
\usepackage{graphicx}
\usepackage{caption}
\usepackage{threeparttable}
\usepackage{float}
\usepackage{algpseudocode}
\usepackage{pdflscape}
\usepackage{enumitem}
\usepackage{xcolor}
\usepackage[nopatch=footnote]{microtype}
\usepackage{cite}
\usepackage{setspace}
\usepackage{lineno}
\usepackage{xurl}
\usepackage{hyperref}

\Urlmuskip=0mu plus 1mu\relax
\setlength{\emergencystretch}{2em}

\hypersetup{%
  colorlinks=true,
  linkcolor=blue!55!black,
  citecolor=blue!55!black,
  urlcolor=blue!55!black
}

\setlength{\parindent}{0pt}
\setlength{\parskip}{0.7\baselineskip}
\setlist{nosep,leftmargin=*}

\newcommand{\AUPRC}{\mathrm{AUPRC}}
\newcommand{\meanSD}[2]{#1\,{\tiny\(\pm\,#2\)}}
\newcommand{\deltaCI}[3]{\shortstack{\ensuremath{#1}\\{\tiny [#2,\,#3]}}}
\newcommand{\ctxmatch}{\shortstack{\ensuremath{\checkmark}\\{\tiny\strut}}}
\newcommand{\ctxexplore}{\shortstack{\ensuremath{\triangle}\\{\tiny\strut}}}
\newcommand{\ctxnomatch}{\shortstack{\ensuremath{\times}\\{\tiny\strut}}}
\newcommand{\ctxmismatch}[1]{\shortstack{\ensuremath{\times}\\{\tiny(#1 duration)}}}
\newcommand{\tableStructLabel}{{\normalsize\textit{Chem.-only}}}
\newcommand{\tableStructBioLabel}{{\normalsize\textit{Chem. + Bio.}}}
\title{Quantifying Biological Complementarity of Perturbation Transcriptomics Beyond Molecular Structure for Toxicity Prediction}
\author{%
  Kyungwhan An\textsuperscript{1,2},
  Ho Gym\textsuperscript{3},
  Eunho Lee\textsuperscript{3},
  Seong Gyu Choi\textsuperscript{3},\\
  Sukyeon Ha\textsuperscript{4},
  Suyeon Kim\textsuperscript{4,*},
  Hwanjo Yu\textsuperscript{1,2,*}\\[0.5em]
  \small\textsuperscript{1}Department of Computer Science and Engineering,\\
  \small Pohang University of Science and Technology (POSTECH), Pohang,
  Republic of Korea\\
  \small\textsuperscript{2}Graduate School of Artificial Intelligence,\\
  \small Pohang University of Science and Technology (POSTECH), Pohang,
  Republic of Korea\\
  \small\textsuperscript{3}BASGENBIO Inc., Seoul, Republic of Korea\\
  \small\textsuperscript{4}School of Computer Science and Engineering,\\
  \small Kyungpook National University, Daegu, Republic of Korea\\[0.5em]
  \small\textbf{Contacts}\\
  \small Kyungwhan An:
  \href{mailto:micolak0115@postech.ac.kr}{micolak0115@postech.ac.kr};
  ORCID: \href{https://orcid.org/0000-0001-9595-5130}{0000-0001-9595-5130}\\
  \small Ho Gym:
  \href{mailto:kkhh4@basgenbio.com}{kkhh4@basgenbio.com};
  ORCID: \href{https://orcid.org/0009-0003-3137-745X}{0009-0003-3137-745X}\\
  \small Eunho Lee:
  \href{mailto:ehlee@basgenbio.com}{ehlee@basgenbio.com};
  ORCID: \href{https://orcid.org/0009-0001-2679-131X}{0009-0001-2679-131X}\\
  \small Seong Gyu Choi:
  \href{mailto:sgchoi@basgenbio.com}{sgchoi@basgenbio.com};
  ORCID: \href{https://orcid.org/0000-0002-2598-2747}{0000-0002-2598-2747} \\
  \small Sukyeon Ha:
  \href{mailto:sukyeon17@knu.ac.kr}
  {sukyeon17@knu.ac.kr};
  ORCID: \href{https://orcid.org/0009-0007-8192-3921}{0009-0007-8192-3921} \\
  \small Suyeon Kim:
  \href{mailto:kimsu55@knu.ac.kr}{kimsu55@knu.ac.kr};
  ORCID: \href{https://orcid.org/0009-0000-9031-4094}{0009-0000-9031-4094};
  Tel.: +82-10-9040-4814\\
  \small Hwanjo Yu:
  \href{mailto:hwanjoyu@postech.ac.kr}{hwanjoyu@postech.ac.kr};
  ORCID: \href{https://orcid.org/0000-0002-7510-0255}{0000-0002-7510-0255};
  Tel.: +82-10-5528-7006\\[0.5em]
  \small\textsuperscript{*}Correspondence:
  \href{mailto:kimsu55@knu.ac.kr}{kimsu55@knu.ac.kr};
  \href{mailto:hwanjoyu@postech.ac.kr}{hwanjoyu@postech.ac.kr}
}
\date{}

\begin{document}
\doublespacing

\maketitle
\linenumbers

\begin{abstract}

Molecular structure-based models provide increasingly competitive toxicity
predictions, yet they do not directly capture cellular responses elicited by
chemical exposure. Perturbation transcriptomics addresses this gap by directly
probing these responses, but its incremental contribution is difficult to
isolate in jointly trained multimodal models. We present a fully interpretable,
model-agnostic framework that freezes an independently developed
molecular toxicity model and uses gene expression responses solely as residual corrections to its prediction. This framework separates biological complementarity
from molecular model retraining or joint optimization, enabling broad
application to future molecular architectures. Across twelve toxicity endpoints and contexts spanning four
cell lines and two exposure times, transcriptomic complementarity was selective
and strongly context dependent. Confirmatory improvements were observed for
SR-MMP and SR-p53, while NR-Aromatase and SR-ARE provided exploratory evidence.
The largest gain was observed for SR-p53 in MCF7 cells at 24 h (\(\Delta\mathrm{AUPRC}=+0.215\)). Biological interpretation of the residual models highlighted endpoint-associated transcriptional programs, including cholesterol homeostasis for NR-Aromatase, oxidative stress for SR-ARE, organelle stress for SR-MMP, and genotoxic stress for SR-p53. At the compound level, transcriptomic residuals rescued assay-active compounds underestimated by the molecular model while also correcting overprediction for assay-inactive compounds. The proposed framework provides a general strategy for quantifying biological complementarity across molecular modeling approaches and a practical basis for prioritizing experimental conditions for future data generation and multimodal toxicity modeling.

\end{abstract}

\noindent\textbf{Keywords:} toxicity prediction; transcriptomics; molecular structure; explainability

\section{Introduction}

Molecular structure is the principal basis for computational toxicity prediction, supported by increasingly expressive representations and learning strategies~\cite{wu2018moleculenet}. Effective prediction depends critically on which structure-property relationships are accessible to the model~\cite{deng2023keyelements}. To capture molecular structure with increasing fidelity and maximize predictive performance across molecular property prediction tasks, molecular representations have progressively evolved from engineered descriptors and fingerprints~\cite{rogers2010extended} to graph~\cite{gilmer2017neural} and tokenized-sequence embeddings~\cite{chithrananda2020chemberta}, geometry-aware~\cite{zhou2023unimol}, image-based~\cite{zeng2022imagemol}, and multiview encodings~\cite{zhang2024mvmrl}. Large-scale pretraining has further improved representation transferability and performance in data-limited toxicity prediction~\cite{mendezlucio2024mole}. Nevertheless, these models remain fundamentally structure-centered and are therefore limited to signals encoded in molecular structure.

Toxicity, however, emerges from the response of an exposed biological system. An initial interaction between a chemical and a molecular target may trigger toxicity, but the resulting adverse effect develops through downstream biological responses across multiple levels of the system~\cite{ankley2010aop}. Molecular structure helps determine physicochemical behavior and the propensity for target engagement, whereas the eventual response is also shaped by dose, exposure duration, biological context, and the functional state of relevant pathways~\cite{niepel2017celltype}. Accordingly, structurally similar compounds can produce different responses across biological contexts, while structurally dissimilar compounds can converge on similar toxic effects~\cite{stumpfe2019activitycliffs}. This creates a fundamental gap for structure-based toxicity prediction. Molecular structure reflects intrinsic toxicological potential but cannot fully determine the response realized under a specific biological context. Cellular response measurements, including pathway-specific assays and perturbation transcriptomics, can therefore provide complementary, context-dependent signals beyond molecular structure.

Motivated by the context dependence of toxicity, prior studies have sought to augment molecular prediction with biological measurements~\cite{liu2023chemicalbiological}, including knowledge-based representations~\cite{park2025cotox}, morphological profiles~\cite{sanchezfernandez2023cloome}, transcriptomic responses~\cite{cattebeke2026integrating}, bioassays~\cite{zhang2025genotoxnet}, and pathological features~\cite{zhong2025hepatotoxicity}. These modalities are typically integrated with molecular features through concatenation or joint optimization. Yet the inclusion of an additional modality neither guarantees that the model will exploit it effectively nor establishes that it contributes predictive value beyond a strong molecular baseline~\cite{hessel2020multimodal}. In chemical-perturbation response prediction, removing or shuffling drug-specific molecular representations has produced little change in performance for unseen compounds, showing that multimodal models can effectively disregard an intended source of signal~\cite{bai2026perturbation}. More strikingly, transcriptomic responses used as pretraining supervision have been reported to impair downstream ADMET prediction~\cite{sypetkowski2024molgps}. These findings shift the focus from increasingly elaborate multimodal integration to a controlled assessment of the incremental predictive value contributed by each modality.

Here, we formulate transcriptomic complementation as a conditional residual problem anchored to a fixed molecular reference. For each endpoint, we first establish a strong molecular model and freeze its predictions before introducing biological data. We then estimate and fix a cohort-specific calibration intercept and learn a ridge-regularized transcriptomic correction on top of the frozen molecular logit. This construction enables a clean paired comparison between molecular-only and transcriptome-complemented predictions on the same test compounds. Because the molecular component is held fixed, any change in prediction arises exclusively from the additive transcriptomic contribution, allowing the incremental predictive value of biological data to be quantified without confounding from model retraining or reweighting. The framework is inherently model-agnostic and enables biological complementarity to be traced from context-level performance to gene- and compound-level contributions.

Applying the framework across twelve toxicity endpoints and perturbation-transcriptomic contexts spanning four cell lines and early and late time points revealed highly selective, context-dependent complementarity. Confirmatory gains were concentrated in SR-MMP and SR-p53, with the largest improvement observed for SR-p53 in MCF7 cells at 24~h (\(\Delta\mathrm{AUPRC}=+0.215\)); NR-Aromatase and SR-ARE provided additional exploratory evidence. Residual models captured endpoint-relevant biology, including cholesterol homeostasis for NR-Aromatase, NRF2-associated oxidative stress for SR-ARE, peroxisomal and mitochondrial stress for SR-MMP, and canonical genotoxic stress for SR-p53. These signals translated into substantial compound-level corrections, including rescue of niclosamide for SR-MMP (\(\Delta p=+0.585\)) and teniposide for SR-p53 (\(\Delta p=+0.703\)), together with downward correction of diphemanil and voriconazole. Transcriptomic complementation also resolved effects not explained by structural similarity alone. Structurally similar teniposide and etoposide diverged, whereas structurally dissimilar niclosamide and gossypol converged toward similar SR-MMP predictions.

\section{Methods}

\subsection{Overview}

\begin{figure}[!htbp]
\centering
\captionsetup{font=footnotesize}
\includegraphics[
  width=0.96\linewidth,
  height=0.42\textheight,
  keepaspectratio
]{Figures/Figure1.eps}
\caption{\textbf{Two-stage framework for context-dependent transcriptomic complementarity.}
Abbreviations: DMSO, dimethyl sulfoxide; trt, treatment.
\textbf{ALT TEXT:}
Flow diagram illustrating molecular structures and Tox21 assays entering a
frozen molecular-prediction stage, and the fixed predictions themselves with LINCS transcriptomic profiles together entering
a context-matched residual-modeling stage, and outputs for context-, gene-,
and compound-level analyses.}
\label{fig:framework-overview}
\end{figure}

We asked whether a measured cellular response improves toxicity prediction
for the same compound beyond what can be predicted from molecular structure
alone (Fig.~\ref{fig:framework-overview}). Stage~1 uses molecular structures
and activity labels to construct an endpoint-specific molecular reference
model, which is subsequently frozen. Stage~2 transfers the frozen molecular
prediction to compound-matched transcriptomic cohorts, estimates a
cohort-level calibration shift, and then learns a regularized transcriptomic
correction while keeping both the molecular prediction and calibration shift
fixed. Molecular-only and transcriptome-complemented predictions are then
compared on the same test compounds to quantify predictive
complementarity. The fitted transcriptomic component is subsequently
interpreted at the gene and compound levels.

\subsection{Input data and biological context construction}

\paragraph{Dataset notation.}

This study uses two input datasets: a molecular dataset with endpoint-specific activity labels \(\mathcal D^{\mathrm{chem}}\), and a transcriptomic dataset indexed by biological context, \(\mathcal D^{\mathrm{bio}}\). For notational simplicity, explicit compound indices are omitted.
Instead, $s$ denotes a structurally identified compound and serves both
as the molecular structure input and as the compound key linking records
across datasets.

The labeled molecular dataset is defined as 

\begin{equation}
\mathcal D^{\mathrm{chem}}
=
\left\{
(s,e,y_e)
\in
\mathcal S\times\mathcal E\times\{0,1\}
:
y_e \text{ is observed for compound } s \text{ at endpoint } e
\right\},
\label{eq:chem-dataset}
\end{equation}

where $\mathcal S$ denotes the set of structurally identified compounds
and $\mathcal E$ denotes the set of twelve binary Tox21 assay tasks.
Each $s\in\mathcal S$ represents a molecular structure together with its
corresponding compound identity, whereas $e\in\mathcal E$ indexes the
assay endpoint. The value $y_e\in\{0,1\}$ denotes the activity label of
compound $s$ for endpoint $e$.

The transcriptomic dataset \(\mathcal D^{\mathrm{bio}}\), whose records do not include activity labels, is defined as

\begin{equation}
\mathcal D^{\mathrm{bio}}
=
\left\{
(s,c,\mathbf x_c)
\in \mathcal S \times \mathcal C \times \mathbb R^{978}
:
\mathbf x_c \text{ is available for } s
\right\}.
\label{eq:bio-dataset}
\end{equation}

where $\mathcal C$ denotes the set of biological contexts, each defined
by a cell background and exposure time. The vector
$\mathbf x_c\in\mathbb R^{978}$ denotes the perturbation-induced
transcriptomic response of compound $s$ in context $c$.

The context-matched cohort is obtained by joining records from
$\mathcal D^{\mathrm{chem}}$ and $\mathcal D^{\mathrm{bio}}$ on the
shared compound $s$ within each biological context $c$:

\begin{equation}
\mathcal D^{\mathrm{chem+bio}}
=
\left\{
(s,e,c,\mathbf x_{ec},y_{ec})
:
(s,e,y_e)\in\mathcal D^{\mathrm{chem}},
\;
(s,c,\mathbf x_c)\in\mathcal D^{\mathrm{bio}}
\right\},
\label{eq:chem+bio-dataset}
\end{equation}

where each matched record contains a molecular structure, an
endpoint-specific activity label, and a transcriptomic response measured
in biological context. For each matched compound, we write
$\mathbf x_{ec}:=\mathbf x_c$ and $y_{ec}:=y_e$ to indicate restriction
to endpoint-context pair $(e,c)$; matching does not alter either
observed quantity.

We further define
\(\mathcal D^{\mathrm{chem-only}}\)
as the scaffold-disjoint subset of \(\mathcal D^{\mathrm{chem}}\) obtained by excluding compounds whose Bemis--Murcko scaffold occurs in \(\mathcal D^{\mathrm{bio}}\). Let $\operatorname{sf}(s)$ denote
the Bemis--Murcko scaffold of compound $s$~\cite{bemis1996properties}. Then

\begin{equation}
\mathcal D^{\mathrm{chem-only}}
=
\left\{
(s,e,y_e)\in\mathcal D^{\mathrm{chem}}
:
\operatorname{sf}(s)\neq\operatorname{sf}(s')
\quad
\forall\,
(s',c,\mathbf x_c)\in\mathcal D^{\mathrm{bio}}
\right\}.
\label{eq:chem-only-dataset}
\end{equation}

The datasets are subsequently restricted by endpoint $e$, context $c$,
or endpoint--context pair $(e,c)$ to define
$\mathcal D^{\mathrm{chem}}_e$,
$\mathcal D^{\mathrm{chem-only}}_e$,
$\mathcal D^{\mathrm{bio}}_c$, and
$\mathcal D^{\mathrm{chem+bio}}_{ec}$, respectively.
The endpoint-specific
$\mathcal D^{\mathrm{chem-only}}_e$ subsets were used exclusively for
Stage~1 molecular-model development. Together,
$\mathcal D^{\mathrm{chem-only}}$ contained 15,695 endpoint-labeled
records corresponding to 1,732 unique compounds
(Table~\ref{tab:stage1-cohort-composition}).

The matched subsets
$\mathcal D^{\mathrm{chem+bio}}_{ec}$ served as the context-matched
Stage~2 cohorts and together contained 67,867 matched records, with the
number of unique compounds varying across endpoint-context pairs
(Tables~\ref{tab:nuclear-receptor-task-cell-counts}
and~\ref{tab:stress-response-task-cell-counts}).

For a cohort containing
$n_{ec}=|\mathcal D^{\mathrm{chem+bio}}_{ec}|$ records, the
978-dimensional transcriptomic response vectors are stacked row-wise to
form the design matrix
$\mathbf X_{ec}\in\mathbb R^{n_{ec}\times978}$, while the corresponding
activity labels form
$\mathbf y_{ec}\in\{0,1\}^{n_{ec}}$.
The frozen Stage~1 model was applied to the molecular structure of each
matched compound to obtain one pre-calibration molecular logit per
compound. These logits were stacked in the same compound order as
$\mathbf X_{ec}$ and $\mathbf y_{ec}$ to form
$\boldsymbol\eta^{\mathrm{chem}}_{ec}\in\mathbb R^{n_{ec}}$.

\paragraph{Small molecules and toxicity assays.}
Twelve Tox21 endpoints were included in the study, comprising seven nuclear-receptor (NR) and five stress-response (SR) endpoints
(Table~\ref{tab:tox21-lincs-configurations})
\cite{huang2016challenge}. The set of compounds
with the available Tox21 label constituted
$\mathcal D^{\mathrm{chem}}_e$ for each endpoint \(e\).
Molecular structures were standardized and assigned Bemis--Murcko
scaffolds using RDKit~v2024.03.2~\cite{landrum2024rdkit}.
To prevent structural overlap with the transcriptomic data, compounds
sharing a scaffold with any compound in
$\mathcal D^{\mathrm{bio}}_c$, for any $c\in\mathcal C$, were excluded,
yielding the scaffold-disjoint endpoint-specific datasets
$\mathcal D^{\mathrm{chem-only}}_e
\subset\mathcal D^{\mathrm{chem}}_e$ defined above.

\paragraph{Perturbation transcriptomics.}
Using LINCS L1000 Level~3 chemical perturbation profiles from HA1E,
HepG2, HT29, and MCF7 cells~\cite{subramanian2017l1000}, we constructed
transcriptional responses to treatment relative to vehicle controls. These
Level~3 profiles, denoted $\mathcal D^{\mathrm{bio}}$, retain normalized
replicate-level expression before differential expression and consensus
transformations, enabling study-specific quality control, replicate
aggregation, and control referencing (Text~S1). All 12,328 genes were retained
during preprocessing; predictive modeling was restricted to the 978 measured
landmark genes.

\paragraph{Context-matched cohort.}
Correspondence between Tox21 and LINCS was evaluated in two stages:
compounds were first matched across datasets, after which each LINCS
profiling context was assessed for biological correspondence with the
associated Tox21 endpoint. Molecular correspondence was established by canonicalizing Tox21 and
LINCS compounds and matching the 14-character InChIKey connectivity
block (Tables~\ref{tab:nuclear-receptor-task-cell-counts}
and~\ref{tab:stress-response-task-cell-counts}).
Restricting each $\mathcal D^{\mathrm{bio}}_c$ to compounds with an
available activity label for endpoint $e$ yielded the context-matched
Stage~2 cohort $\mathcal D^{\mathrm{chem+bio}}_{ec}$. Biological correspondence between Tox21 assays and LINCS contexts was
classified as exact, tissue-proxy, or unmatched
(Table~\ref{tab:tox21-lincs-configurations}).

\clearpage

\begin{landscape}
\begin{table}[p]
\centering
\caption{\textbf{Context correspondence between Tox21 and LINCS datasets.}
The context score summarizes cellular, temporal, and mechanistic correspondence between Tox21 assays and available LINCS profiles; parentheses indicate component scores. Overall context scores of 1-4 are displayed as number of stars ($\bigstar$), while dash (\textit{--}) denotes a score of 0. Exposure times were retained as reported, with the unreported time for SR-ATAD5 left unspecified. \textbf{ALT TEXT:}
Table comparing cellular, temporal, and mechanistic correspondence between
twelve Tox21 assays and available LINCS contexts, with combined context scores
from zero to four.
}
\label{tab:tox21-lincs-configurations}
\fontsize{8.2}{9.6}\selectfont
\setlength{\tabcolsep}{2.8pt}
\renewcommand{\arraystretch}{1.3}
\begin{tabular*}{\linewidth}{@{\extracolsep{\fill}} c l l c l l l c l c @{}}
\toprule
& \multirow[c]{2}{*}{Endpoint} &
\multicolumn{2}{c}{Tox21} &
\multicolumn{2}{c}{LINCS} &
\multicolumn{3}{c}{Context correspondence} &
\multirow[c]{2}{*}{Context score} \\
\cmidrule(lr){3-4}\cmidrule(lr){5-6}\cmidrule(lr){7-9}
& & Cell & Time &
Cell & Time & Cell & Time & Mechanistic & \\
\midrule
\multirow[c]{7}{*}{\rotatebox[origin=c]{90}{Nuclear receptor}}
& AR & MDA-kb2 & 16~h & MCF7 & 6 and 24~h & Breast-Proxy (+1) & Match (+1) & Receptor binding (+0) & $\bigstar\bigstar$ \\
& AR-LBD & HEK293T & 16~h & HA1E & 6 and 24~h & Kidney-Proxy (+1) & Match (+1) & Receptor binding (+0) & $\bigstar\bigstar$ \\
& AhR & HepG2 & 24~h & HepG2 & 6~h & Exact (+2) & No match (+0) & Receptor binding (+0) & $\bigstar\bigstar$ \\
& Aromatase & MCF7 (AroER) & 24~h & MCF7 & 6 and 24~h & Exact (+2) & Match (+1) & Enzymatic activity (+0) & $\bigstar\bigstar\bigstar$ \\
& ER & BG1 & 22~h & \textit{-} & \textit{-} & No match (+0) & No match (+0) & Receptor binding (+0) & \textit{--} \\
& ER-LBD & HEK293 & 18~h & HA1E & 6 and 24~h & Kidney-Proxy (+1) & Match (+1) & Receptor binding (+0) & $\bigstar\bigstar$ \\
& PPAR-\(\gamma\) & HEK293 & 17~h & HA1E & 6 and 24~h & Kidney-Proxy (+1) & Match (+1) & Receptor binding (+0) & $\bigstar\bigstar$ \\
\midrule
\multirow[c]{5}{*}{\rotatebox[origin=c]{90}{Stress response}}
& ARE & HepG2 & 16~h & HepG2 & 6~h & Exact (+2) & No match (+0) & Transcriptional response (+1) & $\bigstar\bigstar\bigstar$ \\
& ATAD5 & HEK293T & \textit{-} & HA1E & 6 and 24~h & Kidney-Proxy (+1) & No match (+0) & Transcriptional response (+1) & $\bigstar\bigstar$ \\
& HSE & HeLa & 5~h & \textit{-} & \textit{-} & No match (+0) & No match (+0) & Transcriptional response (+1) & $\bigstar$ \\
& MMP & HepG2 & 1~h & HepG2 & 6~h & Exact (+2) & No match (+0) & Mitochondrial depolarization (+1) & $\bigstar\bigstar\bigstar$ \\
& p53 & HCT116 & 16~h & HT29 & 6 and 24~h & Colorectal-Proxy (+1) & Match (+1) & Transcriptional response (+1) & $\bigstar\bigstar\bigstar$ \\
\bottomrule
\end{tabular*}
\end{table}
\end{landscape}
\clearpage

\subsection{Model specification}

For each endpoint $e$, the analysis comprised two sequential stages, each
consisting of two steps. Stage~1 constructed an endpoint-specific molecular
reference ensemble $\mathcal M_e^\star$ and applied the frozen ensemble to
each context-matched cohort $\mathcal D^{\mathrm{chem+bio}}_{ec}$, yielding
molecular structure-based predictions $\mathbf p^{\mathrm{chem}}_{ec}$ and
corresponding logits $\boldsymbol\eta^{\mathrm{chem}}_{ec}$.

Stage~2 performed 20 repetitions of five-fold nested cross-validation
(CV) to calibrate the molecular reference and estimate the incremental
predictive contribution of transcriptomic responses. Let $i=1,\ldots,20$
denote the repetition and $j=1,\ldots,5$ the outer fold, with
$u=(i,j)$ indexing the resulting 100 outer-fold evaluations.
Within each $u$, molecular calibration was parameterized by
$\widehat a_{ec}^{(u)}$ and transcriptomic complementation by
$\widehat{\boldsymbol\beta}_{ec}^{(u)}$
(Figure~\ref{fig:nested-validation}). Conceptually,

\begin{equation}
\begin{gathered}
\underbrace{
(\mathcal D^{\mathrm{chem-only}}_e,\mathcal M)
\underset{\mathrm{train}}{\longrightarrow}
\{(\boldsymbol\theta_{em},p^{\mathrm{chem}}_{em})\}_{m\in\mathcal M}
\underset{\mathrm{ensemble}}{\longrightarrow}
\mathcal M_e^\star
}_{\text{Stage 1a: ensemble construction}}
\\[2ex]
%
\underbrace{
(\mathcal D^{\mathrm{chem+bio}}_{ec},\mathcal M_e^\star)
\underset{\mathrm{infer}}{\longrightarrow}
\mathbf p^{\mathrm{chem}}_{ec}
\underset{\mathrm{logit}}{\longrightarrow}
\boldsymbol\eta^{\mathrm{chem}}_{ec}
}_{\text{Stage 1b: molecular structure-based prediction}}
\\[2ex]
%
(\mathcal D^{\mathrm{chem+bio}}_{ec},
\boldsymbol\eta^{\mathrm{chem}}_{ec})
\underset{\mathrm{repeated\ nested\ CV}}{\longrightarrow}
\boxed{
\begin{aligned}
&
\underbrace{
\boldsymbol\eta^{\mathrm{chem},(u)}_{ec}
\overset{\widehat a_{ec}^{(u)}}{\longrightarrow}
\boldsymbol\eta^{\mathrm{chem,cal},(u)}_{ec}
}_{\text{Stage 2a: molecular reference calibration}}
\\[2ex]
&
\underbrace{
(\widetilde{\mathbf X}^{(u)}_{ec},
\boldsymbol\eta^{\mathrm{chem,cal},(u)}_{ec})
\overset{\widehat{\boldsymbol\beta}_{ec}^{(u)}}{\longrightarrow}
\boldsymbol\eta^{\mathrm{bio},(u)}_{ec}
}_{\text{Stage 2b: transcriptomic complementation}}
\end{aligned}
}
\\[-0.5ex]
\qquad
u=(i,j),\quad
i=1,\ldots,20,\quad
j=1,\ldots,5.
\end{gathered}
\label{eq:two-stage-overview}
\end{equation}

\paragraph{Stage 1a: Molecular ensemble construction.}

Stage~1a constructed an endpoint-specific molecular reference from a model
set \(\mathcal M\) spanning complementary molecular
representations and learning architectures. For each model \(m\in\mathcal M\),
the standardized molecular structure \(s\) was mapped to a model-specific
representation \(\phi_m(s)\), comprising a fingerprint vector, a molecular
graph, or a tokenized SMILES sequence.

The model set comprised twelve distinct configurations, including logistic regression and
XGBoost~\cite{chen2016xgboost} applied to MACCS
keys~\cite{durant2002mdl}, Morgan (ECFP4)
fingerprints~\cite{rogers2010extended}, or RDKit topological
fingerprints~\cite{landrum2024rdkit}; Chemprop~\cite{yang2019analyzing};
pretrained graph neural networks (GCN, GAT, GIN, and
GraphSAGE)~\cite{hu2020strategies}; and
ChemBERTa~\cite{chithrananda2020chemberta}. Molecular standardization,
fingerprint generation, and structural processing were performed using
RDKit~v2024.03.2~\cite{landrum2024rdkit}. Model-specific representations,
architectures, and training configurations are provided in
Tables~\ref{tab:stage1-classical-settings}
and~\ref{tab:stage1-deep-settings}.

For endpoint $e\in\mathcal E$, each model $m\in\mathcal M$ was fitted on
the endpoint-labeled training subset of
$\mathcal D^{\mathrm{chem-only}}_e$, with the validation
subset used for model-specific hyperparameter selection and early
stopping where applicable. The fitted model, parameterized by
$\boldsymbol\theta_{em}$, produced the endpoint-toxicity probability

\begin{equation}
p^{\mathrm{chem}}_{em}(s)
=
m\!\left(
\phi_m(s);
\boldsymbol\theta_{em}
\right).
\label{eq:individual-molecular-prediction}
\end{equation}

The twelve fitted models were evaluated on the held-out
ensemble-selection subset and ranked by AUPRC. Let
$m_{e,(1)},\ldots,m_{e,(12)}$ denote this ordering, such that
$\operatorname{AUPRC}(m_{e,(i)})
\geq
\operatorname{AUPRC}(m_{e,(i+1)})$.
For each ensemble size $k=1,\ldots,12$, the top-$k$ models
were combined by unweighted averaging,

\begin{equation}
\mathcal M_e^{(k)}
=
\{m_{e,(1)},\ldots,m_{e,(k)}\},
\qquad
p^{\mathrm{chem},(k)}_e(s)
=
\frac{1}{k}
\sum_{m\in\mathcal M_e^{(k)}}
p^{\mathrm{chem}}_{em}(s).
\label{eq:candidate-molecular-ensemble}
\end{equation}

A common ensemble size was selected by maximizing macro-AUPRC across the
twelve endpoints on the held-out ensemble-selection subset:

\begin{equation}
\overline{\operatorname{AUPRC}}^{(k)}
=
\frac{1}{|\mathcal E|}
\sum_{e\in\mathcal E}
\operatorname{AUPRC}
\!\left(p^{\mathrm{chem},(k)}_e\right),
\qquad
k^\star
=
\operatorname*{arg\,max}_{k\in\{1,\ldots,12\}}
\overline{\operatorname{AUPRC}}^{(k)},
\qquad
\mathcal M_e^\star
=
\mathcal M_e^{(k^\star)}.
\label{eq:ensemble-selection}
\end{equation}

The resulting common size was $k^\star=4$, while the constituent models
remained endpoint specific because model ranking was performed separately
within each endpoint. After ensemble selection, the constituent models,
fitted parameters, and ensemble membership were frozen and transferred
unchanged to Stage~1b.

\paragraph{Stage 1b: Molecular structure-based prediction.}

The frozen ensemble $\mathcal M_e^\star$ was applied without retraining to
each matched endpoint-context cohort
$\mathcal D^{\mathrm{chem+bio}}_{ec}$, yielding
$\mathbf p^{\mathrm{chem}}_{ec}\in(0,1)^{n_{ec}}$.
For additive modeling on the log-odds scale, probabilities were transformed as

\begin{equation}
\boldsymbol\eta^{\mathrm{chem}}_{ec}
=
\operatorname{logit}
\!\left[
\operatorname{clip}_{\varepsilon}
\!\left(
\mathbf p^{\mathrm{chem}}_{ec}
\right)
\right],
\qquad
\varepsilon=10^{-7},
\label{eq:mol-logit}
\end{equation}

where
$\operatorname{clip}_{\varepsilon}(p)
=\min\{1-\varepsilon,\max(\varepsilon,p)\}$
and
$\operatorname{logit}(p)=\log[p/(1-p)]$.
Clipping avoided infinite logits when predicted probabilities were equal to 0 or 1, ensuring stable downstream modeling.

\paragraph{Stage 2a: Molecular structure-based calibration.}

Stage~2 was performed independently for each endpoint-context cohort using
20 repetitions of nested cross-validation. In each repetition, molecular
calibration (Stage~2a) and transcriptomic complementation (Stage~2b) were fitted within the inner and outer training partitions.

A five-fold outer loop was used for held-out performance evaluation, with
$\mathcal D_{ec}^{\mathrm{t}_{\mathrm{out}}}$ denoting the outer-training
set and $\mathcal D_{ec}^{\mathrm{held}}$ its outer-test set. Within each
outer-training set, three-fold inner cross-validation was used for
regularization selection, with
$\mathcal D_{ec}^{\mathrm{t}_{\mathrm{in}}}$ denoting the inner-training
set and $\mathcal D_{ec}^{\mathrm{val}}$ its validation set.
Bemis--Murcko scaffolds were disjoint between training and evaluation
partitions in both loops
(Tables~\ref{tab:outer-validation} and~\ref{tab:inner-validation};
Fig.~\ref{fig:nested-validation}).

For either training level
$r\in\{\mathrm{t}_{\mathrm{in}},\mathrm{t}_{\mathrm{out}}\}$, let
$\mathcal D_{ec}^{r}\subset\mathcal D^{\mathrm{chem+bio}}_{ec}$ contain
$n$ compounds with activity labels
$\mathbf y_{ec}^{r}\in\{0,1\}^{n}\setminus
\{\mathbf 0_n,\mathbf 1_n\}$ and corresponding molecular logits
$\boldsymbol\eta_{ec}^{\mathrm{chem},r}\in\mathbb R^n$.
A scalar calibration intercept was estimated as

\begin{equation}
\widehat a_{ec}^{r}
=
\underset{a\in\mathbb R}{\arg\min}\;
\mathcal L_{\mathrm{BCE}}
\!\left(
\mathbf y_{ec}^{r},
\sigma\!\left(
\boldsymbol\eta_{ec}^{\mathrm{chem},r}
+
a\mathbf 1_n
\right)
\right),
\label{eq:calibration-intercept}
\end{equation}

where
\(
\mathcal L_{\mathrm{BCE}}(\mathbf y,\mathbf p)
=
-\frac{1}{n}
\left[
\mathbf y^{\mathsf T}\log(\mathbf p)
+
(\mathbf 1_n-\mathbf y)^{\mathsf T}
\log(\mathbf 1_n-\mathbf p)
\right],
\)
and $\sigma(t)=[1+\exp(-t)]^{-1}$ is applied element-wise. The calibrated
molecular logits were

\begin{equation}
\boldsymbol\eta_{ec}^{\mathrm{chem,cal},r}
=
\boldsymbol\eta_{ec}^{\mathrm{chem},r}
+
\widehat a_{ec}^{r}\mathbf 1_n.
\label{eq:calibrated-molecular-logit}
\end{equation}

\paragraph{Stage 2b: Transcriptomic complementation.}

For each partition $\mathcal D_{ec}^{r}$, transcriptomic responses were
standardized using statistics estimated from the training
partition $\mathcal D_{ec}^{r'}$:
\begin{equation}
\widetilde{\mathbf x}_{ec}^{r}(s)
=
\bigl(
\mathbf x_c(s)-\boldsymbol\mu_{ec}^{r'}
\bigr)
\oslash
\mathbf d_{ec}^{r'},
\qquad
(r,r')
\in
\{
(\mathrm{t_{in}},\mathrm{t_{in}}),
(\mathrm{val},\mathrm{t_{in}}),
(\mathrm{t_{out}},\mathrm{t_{out}}),
(\mathrm{held},\mathrm{t_{out}})
\}.
\label{eq:standardization}
\end{equation}

Here,
$\boldsymbol\mu_{ec}^{r'},\mathbf d_{ec}^{r'}\in\mathbb R^{978}$
are the training-partition mean and standard-deviation vectors,
$\oslash$ denotes element-wise division, and standard deviations below
$10^{-8}$ were replaced by one. Standardized feature vectors and labels
were stacked row-wise to form $\widetilde{\mathbf X}_{ec}^{r}$ and
$\mathbf y_{ec}^{r}$.

With the calibration intercept $\widehat a_{ec}^{r}$ fixed,
transcriptomic coefficients were estimated by ridge-regularized binary
cross-entropy minimization:

\begin{equation}
\widehat{\boldsymbol\beta}_{ec}^{r}(\lambda)
=
\underset{\boldsymbol\beta\in\mathbb R^{978}}{\arg\min}
\left\{
\mathcal L_{\mathrm{BCE}}
\!\left(
\mathbf y_{ec}^{r},
\sigma\!\left(
\boldsymbol\eta^{\mathrm{chem},r}_{ec}
+\widehat a_{ec}^{r}\mathbf 1_n
+\widetilde{\mathbf X}^{r}_{ec}\boldsymbol\beta
\right)
\right)
+
\frac{\lambda}{2}
\|\boldsymbol\beta\|_2^2
\right\},
\qquad
r\in\{\mathrm{t_{in}},\mathrm{t_{out}}\}.
\label{eq:residual-objective}
\end{equation}

Only $\boldsymbol\beta$ was penalized. Within each inner loop,
$\widehat a_{ec}^{\mathrm{t_{in}}}$ was first estimated and fixed, after
which $\widehat{\boldsymbol\beta}_{ec}^{\mathrm{t_{in}}}(\lambda)$ was
fitted over
$\lambda\in\Lambda\subset[10^{-4},10^{8}]$.

For each valid candidate \(\lambda\), mean validation loss across the three inner folds
was

\begin{equation}
\overline{\mathcal L}^{\mathrm{val}}_{ec}(\lambda)
=
\frac{1}{3}
\sum_{\ell=1}^{3}
\mathcal L^{\mathrm{val},\ell}_{ec}(\lambda),
\label{eq:mean-validation-loss}
\end{equation}

\begin{equation}
\mathcal L^{\mathrm{val},\ell}_{ec}(\lambda)
=
\mathcal L_{\mathrm{BCE}}
\!\left(
\mathbf y^{\mathrm{val},\ell}_{ec},
\sigma\!\left(
\boldsymbol\eta^{\mathrm{chem,val},\ell}_{ec}
+
\widehat a_{ec}^{\mathrm{t_{in}},\ell}
\mathbf 1_{n_{\mathrm{val},\ell}}
+
\widetilde{\mathbf X}^{\mathrm{val},\ell}_{ec}
\widehat{\boldsymbol\beta}_{ec}^{\mathrm{t_{in}},\ell}(\lambda)
\right)
\right),
\label{eq:inner-validation-loss}
\end{equation}

where $\ell$ denotes the inner folds and
$\widetilde{\mathbf X}^{\mathrm{val},\ell}_{ec}$ was standardized using inner-training partition.

The minimum-loss candidate was

\begin{equation}
\lambda^{\min}_{ec}
=
\operatorname*{arg\,min}_{\lambda\in\Lambda}
\overline{\mathcal L}^{\mathrm{val}}_{ec}(\lambda),
\label{eq:lambda-min}
\end{equation}

and the outer-fold regularization parameter was selected by the
one-standard-error (1-SE) rule:

\begin{equation}
\widehat\lambda_{ec}
=
\max
\left\{
\lambda\in\Lambda:
\overline{\mathcal L}^{\mathrm{val}}_{ec}(\lambda)
\le
\overline{\mathcal L}^{\mathrm{val}}_{ec}
\!\left(\lambda^{\min}_{ec}\right)
+
\operatorname{SE}_{ec}
\!\left(\lambda^{\min}_{ec}\right)
\right\}.
\label{eq:lambda-selection}
\end{equation}

The corresponding standard error was

\begin{equation}
\operatorname{SE}_{ec}(\lambda)
=
\sqrt{
\frac{1}{3(3-1)}
\sum_{\ell=1}^{3}
\left[
\mathcal L_{ec}^{\mathrm{val},\ell}(\lambda)
-
\overline{\mathcal L}_{ec}^{\mathrm{val}}(\lambda)
\right]^2
}.
\label{eq:validation-se}
\end{equation}

The 1-SE rule therefore selected the largest $\lambda$ whose mean
validation loss remained within one standard error of the minimum,
favoring stronger regularization to improve stability and reduce overfitting in the
high-dimensional transcriptomic setting.

After selecting $\widehat\lambda_{ec}$, standardization statistics and the
calibration intercept were re-estimated on the complete outer-training
partition $\mathcal D_{ec}^{\mathrm{t_{out}}}$.
With $\widehat a_{ec}^{\mathrm{t_{out}}}$ fixed,
$\boldsymbol\beta$ was refitted using $\widehat\lambda_{ec}$ and applied
to $\mathcal D_{ec}^{\mathrm{held}}$. For each held-out compound $s$,

\begin{equation}
\begin{aligned}
p^{\mathrm{chem,cal}}_{ec}(s)
&=
\sigma\!\left(
\eta^{\mathrm{chem}}_{ec}(s)
+
\widehat a_{ec}^{\mathrm{t_{out}}}
\right),
\\
p^{\mathrm{chem+bio}}_{ec}(s)
&=
\sigma\!\left(
\eta^{\mathrm{chem}}_{ec}(s)
+
\widehat a_{ec}^{\mathrm{t_{out}}}
+
(\widetilde{\mathbf x}_{ec}^{\mathrm{held}}(s))^{\mathsf T}
\widehat{\boldsymbol\beta}_{ec}^{\mathrm{t_{out}}}
(\widehat\lambda_{ec})
\right).
\end{aligned}
\label{eq:final-predictions}
\end{equation}

Because both predictions share the same molecular logit and calibration
intercept, their logit difference equals the transcriptomic contribution:

\begin{equation}
\operatorname{logit}
\!\left(p^{\mathrm{chem+bio}}_{ec}(s)\right)
-
\operatorname{logit}
\!\left(p^{\mathrm{chem,cal}}_{ec}(s)\right)
=
\eta^{\mathrm{bio}}_{ec}(s),
\qquad
\eta^{\mathrm{bio}}_{ec}(s)
=
(\widetilde{\mathbf x}_{ec}^{\mathrm{held}}(s))^{\mathsf T}
\widehat{\boldsymbol\beta}_{ec}^{\mathrm{t_{out}}}
(\widehat\lambda_{ec}).
\label{eq:logit-contribution}
\end{equation}

Across the 20 repetitions, each compound received 20 held-out molecular-only
and transcriptome-complemented predictions, which were retained for
compound-level post-hoc analyses. In the 6-h calibration analysis, calibration
left AUPRC unchanged ($\Delta\mathrm{AUPRC}=0$;
Table~\ref{tab:s1-calibration-selected-endpoints}), as expected because an
intercept shift preserves ranking within each held-out partition.

\subsection{Outputs and post-hoc analyses}

\paragraph{Context-dependent prediction.}

Transcriptomic complementarity was evaluated on outer held-out compounds.
Twenty repetitions of five-fold outer cross-validation yielded 100 paired
evaluations for each endpoint-context pair. AUPRC was used as the
primary metric because of class imbalance~\cite{saito2015precision}.
Complementarity was quantified by \(\Delta\mathrm{AUPRC}\), with positive
values indicating improved ranking after adding the transcriptomic term.
Calibration was assessed separately using
\(\Delta\mathrm{LogLoss}\), with negative values indicating improved probabilistic fit. Paired fold-level effects were
tested using a two-sided Nadeau-Bengio corrected resampled \(t\)-test~\cite{nadeau2003inference}. Benjamini-Hochberg adjustment was applied across
the twelve endpoints separately within each cell line, exposure time, and
analysis cohort.

\paragraph{Gene-level analysis.}

For each endpoint-context pair, the fitted ridge coefficient vectors were
averaged across the 100 outer training fits to obtain
\(\overline{\boldsymbol\beta}_{ec}\in\mathbb R^{978}\). The entries of
\(\overline{\boldsymbol\beta}_{ec}\) summarize the direction and magnitude
of the conditional predictive association of the genes beyond
the molecular reference model.

\paragraph{Compound-level transcriptomic contribution analysis.}

For each outer held-out compound, the total transcriptomic effect was
quantified by the change from the molecular-only to the
transcriptome-complemented probability in
Eq.~\eqref{eq:final-predictions}. Assay-active compounds were classified as rescued
when this probability increased by more than \(0.02\) on average and the
change was positive in all 20 repetitions. Assay-inactive compounds were
classified as corrected when the probability decreased by more than \(0.02\)
on average and the change was negative in all 20 repetitions. Changes were
tested using one-sided Wilcoxon signed-rank tests with Benjamini-Hochberg
correction within each endpoint~\cite{benjamini1995fdr};
compounds with \(q<0.05\) were retained.

For gene-level interpretation, the scalar contribution
\(\eta^{\mathrm{bio}}_{ec}\) was decomposed into the gene-wise contribution
vector
\begin{equation}
\boldsymbol{\psi}^{\mathrm{bio}}_{ec}
=
\widetilde{\mathbf x}_{ec}
\odot
\widehat{\boldsymbol\beta}_{ec}
\in\mathbb R^{978},
\label{eq:gene-wise-transcriptomic-contribution}
\end{equation}
where \(\odot\) denotes element-wise multiplication. The entries sum to the total transcriptomic contribution:
\begin{equation}
\mathbf 1_{978}^{\mathsf T}
\boldsymbol{\psi}^{\mathrm{bio}}_{ec}
=
\eta^{\mathrm{bio}}_{ec}.
\label{eq:sum-gene-contributions}
\end{equation}

\paragraph{Compound-level pathway enrichment analysis.}
For each compound, the gene-wise contribution vectors obtained across its
held-out evaluations were averaged to obtain
\(\overline{\boldsymbol{\psi}}^{\mathrm{bio}}_{ec}\in\mathbb R^{978}\).
Weighted gene-set enrichment analysis (GSEA)~\cite{subramanian2005gsea}
was then applied to the 978 genes ranked by the
entries of
\(\overline{\boldsymbol{\psi}}^{\mathrm{bio}}_{ec}\), using Hallmark gene sets~\cite{liberzon2015hallmark}.
Normalized enrichment scores (NES) were evaluated
against 10,000 size-matched random gene sets. Empirical \(P\)-values were
computed with finite-sampling correction~\cite{phipson2010permutation}
and adjusted using the Benjamini-Hochberg
procedure within each endpoint--pathway combination. Figure~\ref{fig:gene-atlas}
reports confirmatory enrichment at \(q<0.05\) and exploratory enrichment at
\(0.05\leq q<0.10\).

\section{Results}

\subsection{Endpoint-specific ensembles establish the molecular reference}

We first established a molecular reference model against which the incremental
predictive contribution of transcriptomics could be evaluated. Two potential
sources of apparent improvement were addressed before introducing
transcriptomic features: (i) an insufficiently strong molecular model and (ii)
systematic probability-scale differences between the compound-pool of Stage~1 model development and the independent compound-matched cohorts in Stage~2. We therefore constructed endpoint-specific molecular ensembles and then assessed cohort-specific probability shifts after transfer. Together, endpoint-specific ensembling and cohort-specific calibration provided a stable molecular reference model for evaluating transcriptomic complementarity.

No single molecular model performed best across all endpoints, motivating an endpoint-specific ensemble. We compared twelve molecular models across twelve Tox21 endpoints using 1,732 compounds. Morgan-XGBoost achieved the highest nuclear-receptor (NR) macro-AUPRC among individual models at \(0.381\), whereas GIN achieved the highest stress-response (SR) macro-AUPRC at \(0.440\); the best-performing
architecture also differed across individual endpoints (Table~\ref{tab:structure-ensemble-auprc}).

Models were therefore ranked separately within each endpoint and combined by
equal-weight probability averaging. Macro-AUPRC increased from \(0.443\) for
the endpoint-specific best single models (\(k=1\)) to \(0.483\) at \(k=4\),
with corresponding NR and SR macro-AUPRCs of \(0.465\) and \(0.508\),
respectively (Fig.~\ref{fig:ensemble-size-selection}). The four-model ensemble outperformed the highest-ranked
individual model for 10 of 12 endpoints, whereas performance declined as
additional lower-ranked models were included. We therefore selected a common
ensemble size of \(k=4\), while retaining endpoint-specific constituent
models, and froze the resulting ensembles before Stage~2 analysis.

We then examined whether transfer to the compound-matched cohorts introduced cohort-specific calibration shifts. Cohort-level intercept adjustment consistently reduced held-out log loss across all evaluated 6-h endpoint--context pairs (\(\Delta\)log loss: \(-0.2372\) to \(-0.0577\); all \(q<0.05\); Table~\ref{tab:s1-calibration-selected-endpoints}), indicating improved probabilistic calibration. Because the intercept adds a constant to all molecular logits within a cohort, compound rankings were preserved exactly and AUPRC remained unchanged (\(\Delta\mathrm{AUPRC}=0\); Table~\ref{tab:s1-calibration-selected-endpoints}).

\clearpage

\begingroup
\begin{landscape}
\begin{table}[p]
\centering
\caption{\textbf{Molecular-model performance on twelve Tox21 endpoints.}
Endpoint AUPRC and group macro-AUPRC are shown for nuclear-receptor (NR)
and stress-response (SR) endpoints. Bold and underline mark column-best
and second-best values; random-baseline AUPRC equals endpoint activity prevalence. \textbf{ALT TEXT:}
Table comparing AUPRC for twelve molecular-model configurations and a
top-four ensemble across seven nuclear-receptor and five stress-response
toxicity endpoints.}
\label{tab:structure-ensemble-auprc}
\fontsize{9.0}{10.0}\selectfont
\setlength{\tabcolsep}{1.0pt}
\renewcommand{\arraystretch}{1.20}
\resizebox{\linewidth}{!}{%
\begin{tabular}{@{}ccc*{8}{c}*{6}{c}@{}}
\toprule
&
\multirow[c]{2}{*}{Model}
& \multirow[c]{2}{*}{\shortstack{Input\\representation}}
& \multicolumn{8}{c}{Nuclear receptor (NR)}
& \multicolumn{6}{c}{Stress response (SR)} \\
\cmidrule(lr){4-11}\cmidrule(lr){12-17}
& & & AR
& AR-LBD
& AhR
& Aromatase
& ER
& ER-LBD
& PPAR-\(\gamma\)
& Overall
& ARE
& ATAD5
& HSE
& MMP
& p53
& Overall \\
\midrule
& Random baseline
& Held-out activity prevalence
& 0.038 & 0.040 & 0.159 & 0.074 & 0.140 & 0.046 & 0.045 & 0.077
& 0.279 & 0.058 & 0.089 & 0.200 & 0.136 & 0.152 \\

\midrule
\multirow[c]{6}{0.45cm}{\centering\rotatebox[origin=c]{90}{\scriptsize\bfseries Classical ML}}
& \multirow{3}{*}{Logistic regression}
& MACCS keys
& 0.696 & 0.273 & 0.319 & 0.206 & 0.370 & 0.130 & 0.057 & 0.293
& 0.377 & 0.086 & 0.154 & 0.473 & 0.292 & 0.276 \\

& & Morgan fingerprint
& 0.775 & 0.231 & 0.491 & 0.252 & 0.309 & 0.150 & 0.234 & 0.349
& 0.450 & 0.205 & 0.226 & 0.537 & 0.327 & 0.349 \\

& & RDKit topological fingerprint
& 0.706 & 0.339 & 0.445 & 0.251 & 0.238 & 0.194 & 0.193 & 0.338
& 0.470 & 0.177 & 0.144 & 0.552 & 0.286 & 0.326 \\

\addlinespace[6pt]
& \multirow{3}{*}{XGBoost}
& MACCS keys
& 0.616 & 0.374 & 0.375 & 0.360 & 0.384 & \textbf{0.264} & 0.101 & 0.353
& 0.539 & 0.110 & 0.111 & 0.513 & 0.339 & 0.323 \\

& & Morgan fingerprint
& 0.764 & \underline{0.381} & 0.424 & \underline{0.370} & 0.331 & 0.153 & 0.246 & \underline{0.381}
& 0.428 & 0.159 & 0.157 & 0.474 & 0.247 & 0.293 \\

& & RDKit topological fingerprint
& 0.616 & 0.341 & 0.472 & 0.221 & 0.307 & 0.195 & 0.150 & 0.329
& 0.480 & 0.178 & 0.154 & 0.475 & 0.292 & 0.316 \\

\midrule
\multirow[c]{6}{0.45cm}{\centering\rotatebox[origin=c]{90}{\scriptsize\bfseries Deep Learning}}
& GCN
& \multirow[c]{5}{*}{Molecular graph}
& 0.736 & 0.306 & 0.432 & 0.237 & 0.388 & 0.167 & 0.103 & 0.339
& 0.500 & 0.116 & 0.383 & \underline{0.572} & 0.390 & 0.392 \\

& GAT
&
& 0.138 & 0.087 & 0.365 & 0.165 & 0.263 & 0.049 & 0.073 & 0.163
& 0.371 & 0.140 & 0.129 & 0.351 & 0.211 & 0.240 \\

& GraphSAGE
&
& 0.802 & 0.342 & 0.448 & 0.350 & 0.376 & 0.183 & 0.074 & 0.368
& 0.508 & 0.131 & 0.321 & 0.544 & 0.431 & 0.387 \\

& GIN
&
& \underline{0.818} & 0.370 & \underline{0.512} & 0.309 & 0.386 & 0.166 & 0.072 & 0.376
& \underline{0.569} & \underline{0.212} & \textbf{0.456} & 0.525 & \underline{0.436} & \underline{0.440} \\

& Chemprop
&
& 0.734 & 0.232 & 0.373 & 0.251 & 0.335 & 0.121 & \underline{0.308} & 0.336
& 0.432 & 0.188 & 0.313 & 0.533 & 0.397 & 0.373 \\

\addlinespace[6pt]
& ChemBERTa
& SMILES tokens
& 0.744 & 0.367 & 0.450 & 0.283 & \underline{0.415} & 0.161 & 0.044 & 0.352
& 0.504 & 0.180 & 0.203 & 0.562 & 0.301 & 0.350 \\

\midrule
& Ensemble
& Top-four model predictions
& \textbf{0.821} & \textbf{0.415} & \textbf{0.574} & \textbf{0.444}
& \textbf{0.460} & \underline{0.219} & \textbf{0.322} & \textbf{0.465}
& \textbf{0.618} & \textbf{0.341} & \underline{0.432}
& \textbf{0.622} & \textbf{0.527} & \textbf{0.508} \\
\bottomrule
\end{tabular}%
}
\end{table}
\end{landscape}
\endgroup
\clearpage

\begin{figure}[p]
\centering
\captionsetup{font=footnotesize}
\includegraphics[
  width=0.98\linewidth,
  height=0.76\textheight,
  keepaspectratio
]{Figures/Figure2.eps}
\caption{\textbf{Endpoint-specific molecular ensemble selection.}
\textbf{(A)} Macro-AUPRC versus ensemble size \(K\) for all endpoints and
NR and SR subsets.
\textbf{(B,C)} Endpoint AUPRC trajectories for NR and SR, respectively.
Ensemble constituents are endpoint specific and define the frozen molecular
reference. Abbreviations: NR, nuclear receptor; SR, stress response; AUPRC, area under the precision-recall curve. \textbf{ALT TEXT:}
Line graphs showing held-out AUPRC as ensemble size increases from one to
twelve models. Overall, nuclear-receptor, and stress-response macro-AUPRC values
peak near four models, with separate trajectories shown for each of the
twelve endpoints.}
\label{fig:ensemble-size-selection}
\end{figure}
\clearpage

\subsection{Transcriptomic complementarity is selective across biological
contexts}

We next tested whether transcriptomic complementarity depended on biological
context. Assay metadata provided a descriptive assessment of context correspondence (Table~\ref{tab:tox21-lincs-configurations}). Empirically, positive effects were concentrated in NR-Aromatase, SR-ARE, SR-MMP, and SR-p53 endpoints
(Fig.~\ref{fig:biological-context-matching}). Together, these results indicate that context matching identifies settings in which endpoint-relevant transcriptomic responses are more
likely to be observable, but does not guarantee predictive value beyond molecular structure.

\textbf{NR-Aromatase} showed little evidence of transcriptomic
complementarity at 6~h, with \(\Delta\mathrm{AUPRC}\) ranging from
approximately zero in HA1E to \(+0.022\) in MCF7; none of the contexts
reached statistical significance. At 24~h, the assay-aligned MCF7 context
showed the largest improvement in mean AUPRC (\(\Delta\mathrm{AUPRC}=0.128\), 95\% CI: \(0.026\) to \(0.229\); \(P=0.0141\), \(q=0.0564\); Fig.~\ref{fig:biological-context-matching}A),
meeting the exploratory but not confirmatory threshold. The delayed MCF7 gain is
consistent with endocrine signals
emerging downstream of the proximal aromatase-inhibition~\cite{chen2015aromatase}.

\textbf{SR-ARE} exhibited a selective early response (Fig.~\ref{fig:biological-context-matching}B). In HepG2
cells at 6~h, mean AUPRC increased (\(\Delta\mathrm{AUPRC}=0.101\), 95\% CI: \(0.033\) to \(0.170\);
\(P=0.0043\), \(q=0.0511\)), meeting the exploratory threshold. The other 6-h
contexts yielded changes of at most \(0.007\). At 24~h, none of the contexts showed evidence of a significant change. The early predictive signal in HepG2 is
consistent with rapid nuclear factor erythroid 2-related factor 2
(NRF2)-mediated transcriptional adaptation to oxidative
stress~\cite{shukla2012are}.

\textbf{SR-MMP} showed the broadest transcriptomic complementarity
(Fig.~\ref{fig:biological-context-matching}C). At 6~h, AUPRC improved across
all four cell lines, with \(\Delta\mathrm{AUPRC}\) ranging from \(0.059\) in
HA1E to \(0.108\) in MCF7; all four gains were nominally significant
(\(P<0.05\)). At 24~h, substantial improvements persisted in HA1E
(\(\Delta\mathrm{AUPRC}=0.111\), 95\% CI: \(0.030\) to \(0.191\);
\(P=0.0077\), \(q=0.0925\)) and MCF7
(\(\Delta\mathrm{AUPRC}=0.131\), 95\% CI: \(0.052\) to \(0.209\);
\(P=0.0013\), \(q=0.0079\)), whereas HT29 showed no significant gain.
These results indicate that SR-MMP complementarity was not restricted to the assay-aligned HepG2 context, consistent with mitochondrial perturbation eliciting broader metabolic and stress-response programs across cellular backgrounds~\cite{zorova2018mitochondrial}.

\textbf{SR-p53} showed pronounced context dependence
(Fig.~\ref{fig:biological-context-matching}D). At 6~h,
none of the contexts were significant,
including the HT29 context (\(\Delta\mathrm{AUPRC}=0.089\), 95\% CI:
\(-0.007\) to \(0.185\); \(P=0.0680\), \(q=0.4078\)). At 24~h, HT29 again showed
small, non-significant changes, whereas MCF7 produced the largest gain observed
in the study (\(\Delta\mathrm{AUPRC}=0.215\), 95\% CI: \(0.101\) to \(0.328\);
\(P=0.00029\), \(q=0.0035\)). The stronger MCF7 response is biologically plausible given its
wild-type \textit{TP53} background, whereas HT29 carries the transcriptionally
impaired \textit{TP53} R273H hotspot
mutation~\cite{low2019mutantp53}.

By contrast, all remaining NR endpoints other than NR-Aromatase, together with
SR-ATAD5 and SR-HSE, showed no significant transcriptomic gain in any evaluated
context (Table~\ref{tab:main-auprc}). Several of these assays lacked a specified exact LINCS counterpart, but
absence of a match alone did not explain the results. NR-AhR showed no
significant gain despite an exact HepG2 match. Cellular correspondence was
therefore not sufficient to establish predictive complementarity.

\clearpage

\begin{figure}[p]
\centering
\captionsetup{font=footnotesize}
\includegraphics[
  width=\linewidth,
  height=0.72\textheight,
  keepaspectratio
]{Figures/Figure3.eps}
\caption{\textbf{Context-dependent transcriptomic complementarity.}
\textbf{(A-D)} Performance comparison plots across NR-Aromatase, SR-ARE, SR-MMP, and SR-p53 endpoints, respectively. Left: Bar plots displaying mean AUPRC for molecular-only (hatched) and transcriptome-complemented models (no hatch) across biological contexts. Red bars mark the highest AUPRC\@. Right: Bar plots illustrating paired \(\Delta\mathrm{AUPRC}\) with Nadeau-Bengio-corrected 95\% CI\@. Boxed cells are either exact or proxy matched cell lines defined in Table~\ref{tab:tox21-lincs-configurations}. BH adjustment spans twelve endpoints per context and cohort (\(^{\dagger}0.05\leq q<0.10\), exploratory; \(^{*}q<0.05\) and \(^{**}q\leq0.01\), confirmatory; \(q\geq0.10\), non-significant). Abbreviations: NR, nuclear receptor; SR, stress response; ARE, antioxidant response element;
MMP, mitochondrial membrane potential; AUPRC, area under the precision-recall curve. \textbf{ALT TEXT:}
Paired bar graphs on the left comparing
molecular-only and transcriptome-complemented AUPRC and bar graphs on the right for \(\Delta{\mathrm{AUPRC}}\) with confidence intervals for four endpoints (NR-Aromatase, SR-ARE, SR-MMP, and SR-p53) across biological contexts (cell lines and time points).}
\label{fig:biological-context-matching}
\end{figure}
\clearpage

\subsection{Transcriptomic complementation captures endpoint-relevant biology not fully determined by molecular structure}

Across contexts with reproducible predictive gains, transcriptomic residuals
mapped to distinct cellular programs that corroborated the endpoint biology (Fig.~\ref{fig:gene-atlas}).
NR-Aromatase was associated with cholesterol homeostasis indirectly associated with aromatase pathway, SR-ARE with oxidative stress, SR-MMP with mitochondrial and organelle stress, and SR-p53 with canonical genotoxic stress. The residual framework revealed overlooked the cellular programs underlying transcriptomic complementation that were not able to be captured by molecular structure.

\textbf{NR-Aromatase} was associated with cholesterol homeostasis, consistent with the strong negative coefficient of \textit{PSRC1}, a gene implicated in cholesterol regulation (Fig.~\ref{fig:gene-atlas}A)~\cite{luo2022psrc1}. Statin-family compounds showed concordant negative enrichment (NES: \(-2.048\) to \(-1.978\); all \(q<0.05\)). HMG-CoA reductase inhibition perturbs cholesterol and mevalonate metabolism and can alter estrogen-receptor signaling in ER-positive MCF7 cells~\cite{hewitt2003aromatase,simigdala2016cholesterol}, potentially producing downstream responses that overlap with aromatase inhibition.

\textbf{SR-ARE} was characterized by positive coefficients for \textit{HMOX1}, \textit{TXNRD1}, and \textit{GLRX}, consistent with NRF2-regulated antioxidant responses (Fig.~\ref{fig:gene-atlas}B)~\cite{reichard2007nrf2}. ROS-related pathways were positively enriched for compounds including ritanserin and fulvestrant. Fulvestrant has been reported to alter ROS production and antioxidant-enzyme expression in MCF7-derived cells~\cite{badia2016fulvestrant}.

\textbf{SR-MMP} showed enrichment of peroxisomal programs (Fig.~\ref{fig:gene-atlas}C), consistent with the functional coupling between peroxisomes and mitochondria in energy and redox homeostasis~\cite{fransen2017peroxisome}. Parachlorophenol and bicalutamide showed positive and negative peroxisome enrichment, respectively. Androgen-receptor inhibition has also been linked to altered mitochondrial dynamics and peroxisomal ROS signaling~\cite{shiota2024peroxisome}.

\textbf{SR-p53} showed the strongest concordance with canonical DNA-damage responses (Fig.~\ref{fig:gene-atlas}D). \textit{CDKN1A} was the strongest positive coefficient, consistent with p53-mediated cell-cycle arrest and DNA-damage responses~\cite{allen2014p53}. Thiotepa and chlorambucil showed positive enrichment of the p53 pathway, consistent with their established activity as DNA-damaging alkylating agents~\cite{weber2015dnadamaging,krassnig2023chlorambucil}.

Across endpoints, recurrent enrichment of E2F targets, G2/M checkpoint, TNF-\(\alpha\)/NF-\(\kappa\)B signaling, and mTORC1 indicated shared proliferative and stress-response states superimposed on endpoint-specific biology (Fig.~\ref{fig:gene-atlas}E).

\begin{figure}[p]
\centering
\captionsetup{font=footnotesize}
\includegraphics[
  width=\linewidth,
  height=0.72\textheight,
  keepaspectratio
]{Figures/Figure4.eps}
\caption{\textbf{Genes, pathways, and compounds associated with transcriptomic complementarity.}
\textbf{(A-D)} NR-Aromatase, SR-ARE, SR-MMP, and SR-p53 endpoints. Contexts are MCF7 at 24~h except SR-ARE (HepG2, 6~h).
Left: Horizontal bar plots showing top-10 genes with highest absolute ridge coefficients.
Right: Bar plots depicting compound-specific pathway enrichment of top-5 with largest absolute (\(|\mathrm{NES}|\)) for the selected pathway. 
Bar colors indicate the direction of enrichment while significance was denoted by \(^{\dagger}0.05\leq q<0.10\) (exploratory) and \(^{*}q<0.05\) (confirmatory).
\textbf{(E)} Heatmap showing cross-endpoint GSEA screening across compounds. Color intensity gives the
fraction of assay-active compounds with significance \(q<0.10\). Asterisks mark pathways with significant enrichment for at least one compound with confirmatory significance. The black vertical line demarcates common pathways observed in at least three endpoints from those in one or two with \(q<0.05\). Bold denotes selected hallmark pathways with high endpoint relevance.
Abbreviations: NR, nuclear receptor; SR, stress response; ARE, antioxidant response element; MMP, mitochondrial membrane potential. \textbf{ALT TEXT:}
Five-row multi-panel figure with horizontal bar plots on the left showing top positive and negative ridge coefficients and vertical bar plots on the right depicting
compound-level pathway enrichment and a heatmap of pathway enrichment across four endpoints (NR-Aromatase, SR-ARE, SR-MMP, and SR-p53) at the bottom.}
\label{fig:gene-atlas}
\end{figure}
\clearpage

\subsection{Transcriptomic complementation refines molecular structure-based toxicity predictions through realized cellular responses}

Transcriptomic complementation rescued compounds underestimated by the molecular model and corrected compounds that were overestimated across all four endpoints (Fig.~\ref{fig:structural-pathway-atlas}A), with the proportion benefiting ranging from \(6.3\%\) in NR-Aromatase to \(29.7\%\) in SR-ARE\@. Recurrently rescued compounds included methylene blue, digoxin, niclosamide, and ouabain (Fig.~\ref{fig:structural-pathway-atlas}B), whereas nabumetone, bergenin, cimaterol, and diphemanil were repeatedly corrected within their respective contexts (Fig.~\ref{fig:structural-pathway-atlas}C). However, the largest shifts were often endpoint-specific; for example, SSR-69071 showed a strong rescue in SR-ARE, whereas diphemanil showed a prominent correction in SR-MMP\@. This endpoint dependence is consistent with context-specific cellular responses. Structural-pair analysis further demonstrated that molecular similarity alone could not fully explain the resulting probability shifts (Fig.~\ref{fig:structural-case-comparisons}).

Niclosamide showed a strong SR-MMP rescue (\(\Delta p=+0.585\)), consistent with its protonophoric uncoupling and mitochondrial depolarization~\cite{kumar2018niclosamide}, whereas teniposide showed the largest SR-p53 rescue (\(\Delta p=+0.703\)), consistent with topoisomerase-II-mediated DNA damage~\cite{long1992teniposide}. Methylene blue was rescued across all four endpoints and is known to induce mitochondrial redox stress and Nrf2/ARE signaling, with stronger exposure capable of producing oxidative DNA damage~\cite{gureev2019methylene}.

Diphemanil showed the largest SR-MMP decrease (\(\Delta p=-0.090\)); notably, an independent mitochondrial health assay did not identify it as a membrane-potential-decreasing compound~\cite{zink2020ethanol}. Voriconazole decreased in NR-Aromatase (\(\Delta p=-0.047\)) and has been reported to inhibit recombinant human CYP19 only weakly relative to other medicinal azoles~\cite{trosken2004aromatase}.

Importantly, teniposide and etoposide were structurally similar (Morgan Tanimoto similarity \(=0.75\)) and nearly indistinguishable under the molecular model (\(\Delta p_{\mathrm{Chem.}}=+0.018\)), yet diverged markedly after transcriptomic complementation (\(\Delta p_{\mathrm{Chem.+Bio.}}=+0.422\)), consistent with greater intracellular accumulation of teniposide and its higher potency in inducing DNA-break-associated effects despite their shared topoisomerase-II mechanism~\cite{long1992teniposide,long1985dna}. Conversely, niclosamide and gossypol were structurally dissimilar (Morgan Tanimoto similarity \(=0.09\)), yet their SR-MMP probability difference reduced from \(0.041\) to \(0.013\) after transcriptomic complementation, converging at high predicted toxicity. Despite their distinct molecular structures, both converging on shared mitochondrial dysfunction~\cite{kumar2018niclosamide,martinez1988gossypol}.

\clearpage

\begin{figure}[p]
\centering
\captionsetup{font=footnotesize}
\includegraphics[
  width=\linewidth,
  height=0.72\textheight,
  keepaspectratio
]{Figures/Figure5.eps}
\caption{\textbf{Compound-level predictive rescue and correction.}
\textbf{(A)} Scatter plots showing predicted probabilities of molecular-only and transcriptome-complemented models across NR-Aromatase, SR-ARE, SR-MMP, and SR-p53. 
Contexts are MCF7 at 24~h except SR-ARE (HepG2, 6~h). 
Annotated compounds indicate the rescued (dark red) and corrected (blue) with highest probabilities changes. 
The probabilities themselves are highlighted by open red circles.
\textbf{(B,C)} Cross-endpoint rescue and correction heatmaps ranked by cumulative selected change. 
Each heatmap shows the top-20 compounds plus a low ranking compounds with highest endpoint-specific probability changes. 
Abbreviations: NR, nuclear receptor; SR, stress response; ARE, antioxidant response element; MMP, mitochondrial membrane potential. \textbf{ALT TEXT:}
Three-column multi-panel figure with scatter plots on the left comparing molecular-only and transcriptome-complemented
probabilities for four toxicity endpoints, with rescued and corrected
compounds highlighted. Two heatmaps on the right rankings compounds by cumulative
probability increases for rescue and probability decreases for correction.
}
\label{fig:structural-pathway-atlas}
\end{figure}
\clearpage

\section{Discussion}

\textbf{Transcriptomics added predictive value selectively, mainly for stress response.} This pattern is consistent with an AOP framework~\cite{ankley2010aop}, where molecular structure may be especially informative for proximal receptor- or enzyme-mediated events, whereas transcriptomics can capture
downstream oxidative, mitochondrial, genotoxic, and metabolic responses that
emerge only after cellular perturbation. The limited gains across most
nuclear-receptor endpoints therefore argue against transcriptomics as a
uniformly beneficial modality; instead, support its use where the endpoint
is closely coupled to a measurable cellular response state.

\textbf{The usefulness of transcriptomics depended on biological context,
not tissue similarity alone.} NR-Aromatase and SR-ARE broadly followed the expected
cellular and temporal correspondence, whereas SR-MMP and SR-p53 demonstrated
that tissue similarity was insufficient. The broad SR-MMP gains across
cell lines and the pronounced 24-h SR-p53 gain in MCF7 indicate that pathway
competence, response kinetics, and the ability of a cellular system to
propagate downstream stress signals can outweigh naive context
matching. Context selection for transcriptomic profiling should therefore
consider not only tissue correspondence but also whether the relevant
response pathway is intact and observable at the sampled time.

\textbf{These results establish predictive associations in selected cellular
assays rather than causal mechanisms, limiting broader generalizability.}
Tox21 assays measure specific cellular endpoint activity rather than
organism-level adverse outcomes, while LINCS samples a restricted range of
cell lines, doses, and exposure times. The identified response states therefore
may not generalize across tissues, exposure regimens, or \textit{in vivo}
systems. The endpoint-specific linear models also cannot capture nonlinear
interactions or shared response programs spanning multiple endpoint biologies.
Moreover, experimental stress responses do not inherently distinguish
persistent injury from adaptive or compensatory responses, as illustrated by
anti-oxidants observed within the SR-ARE\@ endpoint. Consequently, Tox21 assay
activity should not be interpreted as definitive evidence of adverse toxicity,
because protective adaptation or therapeutically intended stress may
paradoxically produce similar cellular responses.

\textbf{Future studies should rigorously establish when transcriptomic profiling
provides reproducible value beyond molecular structure.}
Denser dose-time sampling, pathway-competent cellular systems, and orthogonal
functional readouts could help distinguish persistent injury from adaptive or
compensatory responses and resolve shared stress programs across endpoints.
In parallel, multitask learning may better capture shared response programs
while preserving endpoint-specific biology, thereby improving prediction.
Selected stress-response assays may also serve as pharmacodynamic readouts of
beneficial adaptation or intended therapeutic stress, warranting further research.

\section{Conclusion}

Perturbation transcriptomics complemented molecular toxicity prediction only
in selected endpoints and biological contexts. Predictive gains were observed
primarily in stress-response endpoints and depended on cellular background and
exposure time. This complementarity reflected endpoint-relevant cellular
biology rather than mere recapitulation of molecular mechanisms. These
findings support context-guided transcriptomic profiling as a targeted strategy
for improving molecular property prediction.

\section*{Key Points}

\begin{itemize}[itemsep=\baselineskip]

\item We developed a systematic framework to quantify the context-specific incremental value of perturbation transcriptomics for molecular structure-based toxicity prediction.

\item Transcriptomic complementarity was selective rather than universal, with gains concentrated in stress-response endpoints and the largest improvement observed for SR-p53 in MCF7 cells at 24 h (\(\Delta\mathrm{AUPRC}=+0.215\)).

\item The framework made transcriptomic complementarity interpretable by connecting context-level performance gains to endpoint-relevant genes, pathways, and compound-specific prediction changes.

\item These findings identify biological context selection as a key design principle for targeted data generation and multimodal toxicity prediction.

\end{itemize}

\section*{Author contributions}

\textbf{Kyungwhan An:} Conceptualization; Formal analysis; Visualization;
Writing---original draft.
\textbf{Ho Gym:} Funding acquisition; Conceptualization; Methodology; Writing---review \& editing.
\textbf{Eunho Lee:} Conceptualization; Methodology; Writing---review \& editing.
\textbf{Seong Gyu Choi:} Conceptualization; Methodology; Writing---review \& editing.
\textbf{Sukyeon Ha:} Methodology; Visualization;
Writing---original draft.
\textbf{Suyeon Kim:} Supervision; Conceptualization; Methodology; Visualization;
Writing---original draft.
\textbf{Hwanjo Yu:} Supervision; Writing---review \& editing.
All authors reviewed and approved the final manuscript.

\section*{Author notes}

\textbf{Competing interests.}
Ho Gym is the founder of BASGENBIO Inc., Seoul, Republic of Korea, and holds
equity interests in the company. Eunho Lee and Seong Gyu Choi are employees of BASGENBIO Inc. All other authors declare no competing interests.

\textbf{Use of LLM-assisted tools.} Large language model (LLM)-based tools were used for language editing, manuscript refinement, and assistance with code writing. These tools did not replace author
judgment in study design, statistical evaluation, or scientific interpretation. All LLM-assisted outputs were critically reviewed, verified, and revised by the authors, who take full responsibility for the accuracy, validity, interpretation, and integrity of the manuscript, analyses, and figures.

\section*{Author biographies}

Kyungwhan An is a doctoral researcher at Pohang University of Science and
Technology (POSTECH), Republic of Korea. His contributions to this study
included conceptualization, formal analysis, visualization, and preparation
of the original manuscript.

Ho Gym is the founder of BASGENBIO Inc., Seoul, Republic of Korea. His
contributions to this study included funding acquisition, conceptualization,
methodological development, and critical review and revision of the manuscript.

Eunho Lee is affiliated with BASGENBIO Inc., Seoul, Republic of Korea. His
contributions to this study included conceptualization, methodological
development, and critical review and revision of the manuscript.

Seong Gyu Choi is affiliated with BASGENBIO Inc., Seoul, Republic of Korea.
His contributions to this study included conceptualization, methodological
development, and critical review and revision of the manuscript.

Sukyeon Ha is a graduate researcher in the School of Computer Science and
Engineering at Kyungpook National University, Republic of Korea. His
contributions included methodological development, visualization, and
preparation of the original manuscript.

Suyeon Kim is a professor in the School of Computer Science and Engineering at
Kyungpook National University, Republic of Korea. Her contributions included
supervision, conceptualization, methodological development, visualization,
and preparation of the manuscript.

Hwanjo Yu is a professor in the Department of Computer Science and Engineering
and the Graduate School of Artificial Intelligence at Pohang University of
Science and Technology (POSTECH), Republic of Korea. His contributions included
supervision and critical review and revision of the manuscript.

\section*{Funding}

This work was supported by the INNOPOLIS Foundation through the 2025 INNOPOLIS R\&D Support Program (AI Global BigTech Fostering Project), funded by the Ministry of Science and ICT, Republic of Korea [grant number RS-2025-25449754].

\section*{Data availability}

The Library of Integrated Network-based Cellular Signatures (LINCS) L1000
Phase~I and Phase~II datasets are publicly available through the Gene
Expression Omnibus (GEO) under accession numbers GSE92742 and GSE70138,
respectively. Tox21 Data Challenge compound-activity data are publicly
available through the Toxicology in the 21st Century (Tox21) Program of the
National Center for Advancing Translational Sciences (NCATS). Scripts used for
data analysis, statistical evaluation, and figure and table generation are
available at \url{https://github.com/micolak0115/BioTox}.

\bibliographystyle{unsrt}
\bibliography{references}

\clearpage
\section*{Supplementary Materials}

% ============================================================
% Supplementary figures
% ============================================================
\begingroup
\renewcommand{\thefigure}{S\arabic{figure}}
\renewcommand{\theHfigure}{supp.\arabic{figure}}
\setcounter{figure}{0}

\begin{center}
\captionsetup{type=figure,font=footnotesize}
\includegraphics[
  width=\linewidth,
  height=0.84\textheight,
  keepaspectratio
]{Figures/FigureS1.eps}
\captionof{figure}{\textbf{Repeated nested validation of transcriptomic complementarity.}
Twenty repetitions of five-fold outer cross-validation yielded 100 paired
held-out evaluations and 20 held-out predictions per compound.
Three-fold inner validation selected the ridge penalty; scaling, calibration,
and transcriptomic coefficients were then refitted on the outer-training
partition. Transcriptomic complementarity was quantified as the logit
difference between transcriptome-complemented and calibrated molecular
predictions. \textbf{ALT TEXT:}
Flow diagram illustrating 20 repetitions of five-fold outer cross-validation
with three-fold inner validation, outer-training refitting, paired held-out
molecular and transcriptome-complemented predictions, and calculation of
the transcriptomic logit contribution.}
\label{fig:nested-validation}
\end{center}

\clearpage

\begin{center}
\captionsetup{type=figure,font=footnotesize}
\includegraphics[
  width=\linewidth,
  height=0.78\textheight,
  keepaspectratio
]{Figures/FigureS2.eps}
\captionof{figure} {\textbf{Comparison of molecular structural similarity with transcriptomic complementarity.}
\textbf{(A-C)} NR-Aromatase, SR-MMP, and SR-p53. Contexts are MCF7 at 24~h. Left: Scatter plots showing predicted probabilities of molecular-only and transcriptome-complemented models across endpoints. 
Middle: A 2D molecular structure display illustrating a structurally conserved pair with divergence in complementation. 
Right: A 2D molecular structure display depicting a structurally distinct pair that are convergent in complementation.
Each pair portrays shared and exclusive substructures while reporting Tanimoto similarity of Morgan fingerprints and absolute between-compound probability differences. \textbf{ALT TEXT:}
Three-row multi-panel figure for three endpoints (NR-Aromatase, SR-MMP, and SR-p53) with scatter plots on the left showing compound-level predictions and 2D structure displays on the right illustrating structurally similar compound pairs with divergent
transcriptomic effects, and structurally dissimilar pairs with convergent
effects, accompanied by molecular substructures and Morgan Tanimoto similarities
displayed.}
\label{fig:structural-case-comparisons}
\end{center}

\endgroup

\clearpage

% ============================================================
% Supplementary tables
% ============================================================
\setcounter{table}{0}
\renewcommand{\thetable}{S\arabic{table}}

% Needed when hyperref is loaded
\makeatletter
\renewcommand{\theHtable}{supp.\arabic{table}}
\makeatother


% ============================================================
% Table S1
% ============================================================
\begin{landscape}
\begin{table}[p]
\centering
\captionsetup{width=\linewidth}

\caption{\textbf{Stage~1 molecular ensemble construction cohorts.}
Endpoint-specific label availability and class composition are shown for
the pooled molecular dataset and its fixed training, validation,
and held-out ensemble-selection partitions.
Labeled denotes compounds with an available endpoint-specific Tox21 label;
Missing denotes compounds without a label for that endpoint. \textbf{ALT TEXT:}
Table showing labeled, active, inactive, and missing compound counts for
each of twelve Tox21 endpoints in the pooled molecular dataset and its
training, validation, and ensemble-selection partitions.}
\label{tab:stage1-cohort-composition}

\fontsize{9}{10}\selectfont
\setlength{\tabcolsep}{2.2pt}
\renewcommand{\arraystretch}{1.45}
\setlength{\aboverulesep}{1pt}
\setlength{\belowrulesep}{1pt}
\setlength{\arrayrulewidth}{0.6pt}

\begin{tabularx}{\linewidth}{
@{}
>{\centering\arraybackslash}p{0.55cm}
>{\raggedright\arraybackslash}p{1.55cm}|
*{4}{>{\centering\arraybackslash}X}|
*{4}{>{\centering\arraybackslash}X}|
*{4}{>{\centering\arraybackslash}X}|
*{4}{>{\centering\arraybackslash}X}
@{}
}

\toprule

&
\multirow[c]{2}{*}{\textbf{Endpoint}} &
\multicolumn{4}{c|}{\textbf{Molecular pool} $(n=1{,}732)$} &
\multicolumn{4}{c|}{\textbf{Training} $(n=1{,}039)$} &
\multicolumn{4}{c|}{\textbf{Validation} $(n=346)$} &
\multicolumn{4}{c}{\textbf{Ensemble selection} $(n=347)$} \\

\cmidrule(lr){3-6}
\cmidrule(lr){7-10}
\cmidrule(lr){11-14}
\cmidrule(lr){15-18}

& &
Labeled & Active & Inactive & Missing &
Labeled & Active & Inactive & Missing &
Labeled & Active & Inactive & Missing &
Labeled & Active & Inactive & Missing \\

\midrule

\multirow[c]{7}{*}{%
\rotatebox[origin=c]{90}{\bfseries Nuclear receptor}}
& AR
& 1,555 & 107 & 1,448 & 177
& 927 & 83 & 844 & 112
& 316 & 12 & 304 & 30
& 312 & 12 & 300 & 35 \\

& AR-LBD
& 1,393 & 92 & 1,301 & 339
& 838 & 70 & 768 & 201
& 281 & 11 & 270 & 65
& 274 & 11 & 263 & 73 \\

& AhR
& 1,352 & 217 & 1,135 & 380
& 809 & 125 & 684 & 230
& 272 & 49 & 223 & 74
& 271 & 43 & 228 & 76 \\

& Aromatase
& 1,081 & 102 & 979 & 651
& 645 & 65 & 580 & 394
& 232 & 22 & 210 & 114
& 204 & 15 & 189 & 143 \\

& ER
& 1,199 & 194 & 1,005 & 533
& 726 & 132 & 594 & 313
& 237 & 29 & 208 & 109
& 236 & 33 & 203 & 111 \\

& ER-LBD
& 1,434 & 90 & 1,344 & 298
& 866 & 63 & 803 & 173
& 284 & 14 & 270 & 62
& 284 & 13 & 271 & 63 \\

& PPAR-$\gamma$
& 1,268 & 51 & 1,217 & 464
& 766 & 29 & 737 & 273
& 255 & 11 & 244 & 91
& 247 & 11 & 236 & 100 \\

\midrule

\multirow[c]{5}{*}{%
\rotatebox[origin=c]{90}{\bfseries Stress response}}
& ARE
& 1,109 & 269 & 840 & 623
& 684 & 157 & 527 & 355
& 210 & 52 & 158 & 136
& 215 & 60 & 155 & 132 \\

& ATAD5
& 1,478 & 73 & 1,405 & 254
& 885 & 40 & 845 & 154
& 298 & 16 & 282 & 48
& 295 & 17 & 278 & 52 \\

& HSE
& 1,280 & 94 & 1,186 & 452
& 781 & 51 & 730 & 258
& 240 & 20 & 220 & 106
& 259 & 23 & 236 & 88 \\

& MMP
& 1,155 & 262 & 893 & 577
& 695 & 163 & 532 & 344
& 230 & 53 & 177 & 116
& 230 & 46 & 184 & 117 \\

& p53
& 1,391 & 136 & 1,255 & 341
& 836 & 70 & 766 & 203
& 282 & 29 & 253 & 64
& 273 & 37 & 236 & 74 \\

\bottomrule

\end{tabularx}
\end{table}
\end{landscape}
\clearpage

% ============================================================
% Table S2
% ============================================================
\begin{landscape}
\begingroup
\begin{table}[p]
\centering
\captionsetup{width=\linewidth,font=footnotesize}
\caption{\textbf{Stage~2 context-matched cohorts for nuclear-receptor}
\(n\): labeled compounds, the denominator for active percentages.
\(\checkmark\): confirmatory-eligible; \(\triangle\): exploratory;
\(\times\): excluded. Bold active counts are \(<20\); eligibility is not
statistical significance. Counts use the full 6-h cohort and compounds
profiled at both times for 24~h. Total profiled (6~h/24~h):
HA1E, 1,012/861; HepG2, 695/46; HT29, 1,005/606; MCF7, 1,332/1,316. \textbf{ALT TEXT:}
Table showing labeled-compound counts, active counts, activity percentages,
and analysis eligibility for seven nuclear-receptor endpoints across four
cell lines at 6 and 24 hours.}
\label{tab:nuclear-receptor-task-cell-counts}
\fontsize{10.2}{11.2}\selectfont
\setlength{\tabcolsep}{3.5pt}
\renewcommand{\arraystretch}{0.94}
\setlength{\aboverulesep}{1pt}
\setlength{\belowrulesep}{1pt}
\begin{tabular*}{\linewidth}{@{\extracolsep{\fill}}ll|rrrcrrrc@{}} % chktex 44
\toprule
\multirow[c]{2}{*}{\textbf{Endpoint}} &
\multirow[c]{2}{*}{\textbf{Cell}} &
\multicolumn{4}{c}{\textbf{6~h}} &
\multicolumn{4}{c}{\textbf{24~h}} \\
\cmidrule(lr){3-6}\cmidrule(lr){7-10}
& & \(n\) & Active & Active (\%) & Status &
\(n\) & Active & Active (\%) & Status \\
\midrule
\multirow[c]{4}{*}{AR} & HA1E & 968 & 72 & 7.44 & \(\checkmark\) & 830 & 59 & 7.11 & \(\checkmark\) \\
 & HepG2 & 668 & 46 & 6.89 & \(\checkmark\) & 45 & \textbf{5} & 11.11 & \(\times\) \\
 & HT29 & 960 & 72 & 7.50 & \(\checkmark\) & 585 & 44 & 7.52 & \(\checkmark\) \\
 & MCF7 & 1259 & 82 & 6.51 & \(\checkmark\) & 1243 & 78 & 6.28 & \(\checkmark\) \\
\midrule
\multirow[c]{4}{*}{AR-LBD} & HA1E & 869 & 49 & 5.64 & \(\checkmark\) & 747 & 39 & 5.22 & \(\checkmark\) \\
 & HepG2 & 598 & 37 & 6.19 & \(\checkmark\) & 31 & \textbf{1} & 3.23 & \(\times\) \\
 & HT29 & 862 & 47 & 5.45 & \(\checkmark\) & 523 & \textbf{19} & 3.63 & \(\triangle\) \\
 & MCF7 & 1145 & 53 & 4.63 & \(\checkmark\) & 1129 & 51 & 4.52 & \(\checkmark\) \\
\midrule
\multirow[c]{4}{*}{AhR} & HA1E & 849 & 102 & 12.01 & \(\checkmark\) & 727 & 91 & 12.52 & \(\checkmark\) \\
 & HepG2 & 585 & 65 & 11.11 & \(\checkmark\) & 36 & \textbf{6} & 16.67 & \(\times\) \\
 & HT29 & 841 & 95 & 11.30 & \(\checkmark\) & 516 & 63 & 12.21 & \(\checkmark\) \\
 & MCF7 & 1111 & 112 & 10.08 & \(\checkmark\) & 1099 & 111 & 10.10 & \(\checkmark\) \\
\midrule
\multirow[c]{4}{*}{Aromatase} & HA1E & 759 & 63 & 8.30 & \(\checkmark\) & 655 & 57 & 8.70 & \(\checkmark\) \\
 & HepG2 & 512 & 47 & 9.18 & \(\checkmark\) & 26 & \textbf{14} & 53.85 & \(\times\) \\
 & HT29 & 758 & 63 & 8.31 & \(\checkmark\) & 475 & 45 & 9.47 & \(\checkmark\) \\
 & MCF7 & 1021 & 76 & 7.44 & \(\checkmark\) & 1012 & 76 & 7.51 & \(\checkmark\) \\
\midrule
\multirow[c]{4}{*}{ER} & HA1E & 805 & 132 & 16.40 & \(\checkmark\) & 689 & 112 & 16.26 & \(\checkmark\) \\
 & HepG2 & 540 & 92 & 17.04 & \(\checkmark\) & 23 & \textbf{11} & 47.83 & \(\times\) \\
 & HT29 & 799 & 128 & 16.02 & \(\checkmark\) & 487 & 76 & 15.61 & \(\checkmark\) \\
 & MCF7 & 1067 & 155 & 14.53 & \(\checkmark\) & 1054 & 154 & 14.61 & \(\checkmark\) \\
\midrule
\multirow[c]{4}{*}{ER-LBD} & HA1E & 907 & 55 & 6.06 & \(\checkmark\) & 772 & 40 & 5.18 & \(\checkmark\) \\
 & HepG2 & 620 & 39 & 6.29 & \(\checkmark\) & 39 & \textbf{8} & 20.51 & \(\times\) \\
 & HT29 & 902 & 53 & 5.88 & \(\checkmark\) & 547 & 25 & 4.57 & \(\checkmark\) \\
 & MCF7 & 1194 & 61 & 5.11 & \(\checkmark\) & 1178 & 59 & 5.01 & \(\checkmark\) \\
\midrule
\multirow[c]{4}{*}{PPAR-\(\gamma\)} & HA1E & 811 & 34 & 4.19 & \(\checkmark\) & 696 & 29 & 4.17 & \(\checkmark\) \\
 & HepG2 & 552 & \textbf{19} & 3.44 & \(\triangle\) & 29 & \textbf{2} & 6.90 & \(\times\) \\
 & HT29 & 803 & 29 & 3.61 & \(\checkmark\) & 483 & \textbf{17} & 3.52 & \(\triangle\) \\
 & MCF7 & 1073 & 34 & 3.17 & \(\checkmark\) & 1060 & 33 & 3.11 & \(\checkmark\) \\
\bottomrule
\end{tabular*}
\end{table}
\endgroup
\end{landscape}
\clearpage

% ============================================================
% Table S3
% ============================================================
\begin{landscape}
\begingroup
\begin{table}[p]
\centering
\captionsetup{width=\linewidth}
\caption{\textbf{Stage~2 context-matched cohorts for stress-response}
\(n\): labeled compounds, the denominator for active percentages.
\(\checkmark\): confirmatory-eligible; \(\triangle\): exploratory;
\(\times\): excluded. Bold active counts are \(<20\); eligibility is not
statistical significance. Counts use the full 6-h cohort and compounds
profiled at both times for 24~h. Total profiled (6~h/24~h):
HA1E, 1,012/861; HepG2, 695/46; HT29, 1,005/606; MCF7, 1,332/1,316. \textbf{ALT TEXT:}
Table showing labeled-compound counts, active counts, activity percentages,
and analysis eligibility for five stress-response endpoints across four
cell lines at 6 and 24 hours.}
\label{tab:stress-response-task-cell-counts}
\fontsize{13}{15}\selectfont
\setlength{\tabcolsep}{4pt}
\renewcommand{\arraystretch}{1.0}
\setlength{\aboverulesep}{1pt}
\setlength{\belowrulesep}{1pt}
\begin{tabular*}{\linewidth}{@{\extracolsep{\fill}}ll|rrrcrrrc@{}} % chktex 44
\toprule
\multirow[c]{2}{*}{\textbf{Endpoint}} &
\multirow[c]{2}{*}{\textbf{Cell}} &
\multicolumn{4}{c}{\textbf{6~h}} &
\multicolumn{4}{c}{\textbf{24~h}} \\
\cmidrule(lr){3-6}\cmidrule(lr){7-10}
& & \(n\) & Active & Active (\%) & Status &
\(n\) & Active & Active (\%) & Status \\
\midrule
\multirow[c]{4}{*}{ARE} & HA1E & 690 & 136 & 19.71 & \(\checkmark\) & 584 & 118 & 20.21 & \(\checkmark\) \\
 & HepG2 & 451 & 98 & 21.73 & \(\checkmark\) & 26 & \textbf{15} & 57.69 & \(\times\) \\
 & HT29 & 682 & 129 & 18.91 & \(\checkmark\) & 416 & 80 & 19.23 & \(\checkmark\) \\
 & MCF7 & 923 & 160 & 17.33 & \(\checkmark\) & 913 & 159 & 17.42 & \(\checkmark\) \\
\midrule
\multirow[c]{4}{*}{ATAD5} & HA1E & 911 & 50 & 5.49 & \(\checkmark\) & 778 & 43 & 5.53 & \(\checkmark\) \\
 & HepG2 & 619 & 33 & 5.33 & \(\checkmark\) & 33 & \textbf{6} & 18.18 & \(\times\) \\
 & HT29 & 903 & 50 & 5.54 & \(\checkmark\) & 544 & 30 & 5.51 & \(\checkmark\) \\
 & MCF7 & 1196 & 57 & 4.77 & \(\checkmark\) & 1180 & 55 & 4.66 & \(\checkmark\) \\
\midrule
\multirow[c]{4}{*}{HSE} & HA1E & 814 & 52 & 6.39 & \(\checkmark\) & 697 & 43 & 6.17 & \(\checkmark\) \\
 & HepG2 & 547 & 30 & 5.48 & \(\checkmark\) & 36 & \textbf{6} & 16.67 & \(\times\) \\
 & HT29 & 806 & 50 & 6.20 & \(\checkmark\) & 498 & 31 & 6.22 & \(\checkmark\) \\
 & MCF7 & 1077 & 64 & 5.94 & \(\checkmark\) & 1065 & 62 & 5.82 & \(\checkmark\) \\
\midrule
\multirow[c]{4}{*}{MMP} & HA1E & 723 & 129 & 17.84 & \(\checkmark\) & 613 & 111 & 18.11 & \(\checkmark\) \\
 & HepG2 & 495 & 102 & 20.61 & \(\checkmark\) & 29 & \textbf{19} & 65.52 & \(\times\) \\
 & HT29 & 717 & 121 & 16.88 & \(\checkmark\) & 430 & 61 & 14.19 & \(\checkmark\) \\
 & MCF7 & 958 & 143 & 14.93 & \(\checkmark\) & 947 & 141 & 14.89 & \(\checkmark\) \\
\midrule
\multirow[c]{4}{*}{p53} & HA1E & 884 & 89 & 10.07 & \(\checkmark\) & 757 & 76 & 10.04 & \(\checkmark\) \\
 & HepG2 & 603 & 64 & 10.61 & \(\checkmark\) & 38 & \textbf{15} & 39.47 & \(\times\) \\
 & HT29 & 877 & 85 & 9.69 & \(\checkmark\) & 533 & 57 & 10.69 & \(\checkmark\) \\
 & MCF7 & 1158 & 100 & 8.64 & \(\checkmark\) & 1142 & 98 & 8.58 & \(\checkmark\) \\
\bottomrule
\end{tabular*}
\end{table}
\endgroup
\end{landscape}
\clearpage

% ============================================================
% Table S4
% ============================================================
\begin{table}[p]
\centering
\captionsetup{width=\linewidth}
\caption{\textbf{Classical molecular-model settings.}
Logistic regression and XGBoost were each evaluated independently with
MACCS keys, Morgan (ECFP4) fingerprints, and RDKit topological
fingerprints for every endpoint. Morgan fingerprints used radius 2.
Braces denote hyperparameter search spaces. \textbf{ALT TEXT:}
Table listing input fingerprints, model architecture, hyperparameter
ranges, optimization settings, and training configurations for logistic
regression and XGBoost molecular models.}
\label{tab:stage1-classical-settings}

\footnotesize
\setlength{\tabcolsep}{4pt}
\renewcommand{\arraystretch}{1.05}
\setlength{\arrayrulewidth}{0.6pt}

\begin{tabularx}{\linewidth}{
@{}
>{\raggedright\arraybackslash}p{0.22\linewidth}|
>{\raggedright\arraybackslash}p{0.31\linewidth}
>{\raggedright\arraybackslash}X
@{}
}
\toprule
\textbf{Model family}
& \textbf{Setting}
& \textbf{Specification} \\
\midrule

\textbf{Logistic regression}
& Input representation (dimension)
& MACCS keys (167 bits) \\

&
& Morgan fingerprint, ECFP4 (2,048 bits) \\

&
& RDKit topological fingerprint (2,048 bits) \\

& Architecture
& Linear logistic classifier \\

& Inverse regularization, \(C\)
& \(\{0.001, 0.01, 0.1, 1, 10, 100\}\) \\

& Training budget
& Maximum 2,000 iterations \\

& Optimizer
& L-BFGS \\

& Learning rate
& N/A \\

& Loss function
& L$_2$-regularized binary cross-entropy \\

\midrule

\textbf{XGBoost}
& Input representation (dimension)
& MACCS keys (167 bits) \\

&
& Morgan fingerprint, ECFP4 (2,048 bits) \\

&
& RDKit topological fingerprint (2,048 bits) \\

& Architecture
& Gradient-boosted decision trees (gbtree) \\

& Maximum tree depth
& \(\{3,4,6\}\) \\

& Row / column subsampling
& 0.8 / 0.8 \\

& Training budget
& Maximum 300 rounds / 20-round patience \\

& Optimizer
& Gradient tree boosting \\

& Learning rate
& \(\{0.05,0.1\}\) \\

& Loss function
& Binary cross-entropy (\texttt{binary:logistic}) \\

\bottomrule
\end{tabularx}
\end{table}
\clearpage

% ============================================================
% Table S5
% ============================================================
\begin{table}[p]
\centering
\captionsetup{width=\linewidth,font=footnotesize}
\caption{\textbf{Deep molecular-model settings.}
Chemprop used its multi-hot graph featurizer: 72 atom features (version 2)
encode atomic number, degree, formal charge, chirality, hydrogen count,
hybridization, aromaticity, and atomic mass, whereas 14 bond features encode
bond type, conjugation, ring membership, stereochemistry, and null-bond
status. Pretrained GNNs used atomic number and chirality as categorical
node attributes and bond type and bond direction as categorical edge
attributes, preserving compatibility with the
supervised\_contextpred checkpoint; categorical indices were
mapped to learned embeddings before message passing. \textbf{ALT TEXT:}
Table listing molecular representations, architectures, hidden dimensions,
training schedules, optimization settings, and loss functions for
Chemprop, four pretrained graph neural networks, and ChemBERTa.
}
\label{tab:stage1-deep-settings}

\fontsize{7.5}{8.2}\selectfont
\setlength{\tabcolsep}{3pt}
\renewcommand{\arraystretch}{0.92}
\setlength{\arrayrulewidth}{0.6pt}

\begin{tabularx}{\linewidth}{
@{}
>{\raggedright\arraybackslash}p{0.20\linewidth}|
>{\raggedright\arraybackslash}p{0.31\linewidth}
>{\raggedright\arraybackslash}X
@{}
}
\toprule
\textbf{Model family}
& \textbf{Setting}
& \textbf{Specification} \\
\midrule

\textbf{Chemprop}
& Input representation (dimension)
& Molecular graph; node attributes: multi-hot atom featurizer v2 (72);
edge attributes: multi-hot bond featurizer (14) \\

& Architecture
& Directed message-passing neural network (D-MPNN) \\

& Message-passing depth
& 3 \\

& Hidden dimension
& 300 \\

& Graph readout
& Mean pooling \\

& Dropout
& 0 \\

& Batch size
& 32 \\

& Training budget
& 100 epochs / 20-epoch patience \\

& Optimizer
& Adam \\

& Learning rate
& Noam: \(10^{-4} \rightarrow 10^{-3} \rightarrow 10^{-4}\);
2-epoch warmup \\

& Loss function
& Binary cross-entropy \\

\midrule

\textbf{Pretrained GNNs}
& Input representation (dimension)
& Molecular graph; node attributes: atomic number and chirality (2);
edge attributes: bond type and bond direction (2) \\

& Architecture
& GCN, GAT, GIN, and GraphSAGE \\

& Pretrained checkpoint
& \texttt{supervised\_contextpred} \\

& Message-passing depth
& 5 layers \\

& Hidden dimension
& 300 \\

& Graph readout
& Mean pooling \\

& Dropout
& 0.5 \\

& Batch size
& 32 \\

& Training budget
& 100 epochs / 20-epoch patience \\

& Optimizer
& AdamW \\

& Learning rate
& \(10^{-4}\) \\

& Loss function
& Binary cross-entropy \\

\midrule

\textbf{ChemBERTa}
& Input representation (dimension)
& Tokenized SMILES (maximum 128 tokens) \\

& Architecture
& Transformer-based molecular language model (MLM) \\

& Pretrained checkpoint
& ChemBERTa-100M-MLM \\

& Weight decay
& 0 \\

& Linear warmup
& 6\% of training steps \\

& Batch size
& 32 \\

& Training budget
& 20 epochs / 5-epoch patience \\

& Optimizer
& AdamW \\

& Learning rate
& \(10^{-4}\) \\

& Loss function
& Binary cross-entropy \\

\bottomrule
\end{tabularx}
\end{table}
\clearpage

% ============================================================
% Table S6
% ============================================================
\begin{landscape}
\begingroup

\begin{table}[p]
\captionsetup{font=small,width=\linewidth,skip=6pt}
\caption{\textbf{Repeated outer-loop evaluation.}
Twenty repetitions of five-fold outer cross-validation define 100
outer-fold evaluations for each endpoint--context cohort. Molecular-only
and transcriptome-complemented predictions are evaluated on the same
held-out compounds, with all fitted quantities estimated exclusively
from the corresponding outer-training partition. \textbf{ALT TEXT:}
Algorithm table showing 20 repetitions of five-fold outer
cross-validation, including scaffold-disjoint splitting, inner-loop ridge
selection, outer-training calibration and transcriptomic fitting, and
paired held-out performance evaluation.}
\label{tab:outer-validation}

\fontsize{11}{12}\selectfont
\setlength{\parskip}{0pt}
\setlength{\abovedisplayskip}{3pt}
\setlength{\belowdisplayskip}{3pt}

\noindent\rule{\linewidth}{0.8pt}\par\nobreak

\begin{algorithmic}[1]

\Require Context-matched cohort
$\mathcal D^{\mathrm{chem+bio}}_{ec}$,
frozen molecular logits
$\boldsymbol\eta^{\mathrm{chem}}_{ec}$,
and ridge grid $\Lambda$.

\Ensure Paired held-out molecular-only and transcriptome-complemented
predictions, performance differences, and fitted parameters.

\For{each repetition $i=1,\ldots,20$}

    \State Generate a distinct five-fold scaffold-disjoint outer split.

    \For{each outer fold $j=1,\ldots,5$}

        \State Define fold $j$ as
        $\mathcal D_{ec}^{\mathrm{held}}$
        and the remaining four folds as
        $\mathcal D_{ec}^{\mathrm{t_{out}}}$.

        \State
        $\widehat\lambda_{ec}
        \gets
        \Call{SelectLambda}
        {\mathcal D_{ec}^{\mathrm{t_{out}}},\Lambda}$
        using the three-fold inner procedure
        (Table~\ref{tab:inner-validation}).

        \If{ridge selection fails}
            \State Skip the current outer fold.
        \EndIf

        \State Fit the calibration intercept
        $\widehat a_{ec}^{\mathrm{t_{out}}}$
        on $\mathcal D_{ec}^{\mathrm{t_{out}}}$,
        with the molecular logit coefficient fixed to one
        (Eq.~\ref{eq:calibration-intercept}).

        \State Estimate
        $\boldsymbol\mu_{ec}^{\mathrm{t_{out}}}$ and
        $\mathbf d_{ec}^{\mathrm{t_{out}}}$
        from $\mathcal D_{ec}^{\mathrm{t_{out}}}$,
        and apply them unchanged to
        $\mathcal D_{ec}^{\mathrm{held}}$
        (Eq.~\ref{eq:standardization}).

        \State With
        $\widehat a_{ec}^{\mathrm{t_{out}}}$ fixed, fit
        $\widehat{\boldsymbol\beta}_{ec}^{\mathrm{t_{out}}}
        (\widehat\lambda_{ec})$
        on $\mathcal D_{ec}^{\mathrm{t_{out}}}$
        (Eq.~\ref{eq:residual-objective}).

        \State Compute paired held-out predictions
        $p^{\mathrm{chem,cal}}_{ec}$ and
        $p^{\mathrm{chem+bio}}_{ec}$
        on $\mathcal D_{ec}^{\mathrm{held}}$
        (Eq.~\ref{eq:final-predictions}).

        \State Compute paired held-out
        $\Delta\mathrm{AUPRC}$ and
        $\Delta\mathrm{LogLoss}$, and retain predictions,
        outcomes, and fitted parameters for downstream analyses.

    \EndFor

\EndFor

\State \Return Paired held-out results across all valid outer-fold evaluations.

\end{algorithmic}

\par\nobreak\noindent\rule{\linewidth}{0.8pt}

\end{table}
\endgroup
\end{landscape}
\clearpage


% ============================================================
% Table S7
% ============================================================
\begin{landscape}
\begingroup

\begin{table}[p]
\captionsetup{font=small,width=\linewidth,skip=6pt}
\caption{\textbf{Inner-loop ridge selection.}
Within each outer-training partition, three-fold inner cross-validation
selects the largest ridge penalty whose mean validation BCE lies within
one standard error of the minimum. Standardization, calibration, and
transcriptomic coefficients are estimated exclusively from the
corresponding inner-training partition. Abbreviations: BCE, binary
cross-entropy. \textbf{ALT TEXT:}
Algorithm table showing three-fold inner cross-validation in which
calibration and transcriptomic coefficients are fitted for each ridge
penalty, validation losses are averaged, and the final penalty is selected
using the one-standard-error rule.}
\label{tab:inner-validation}

\fontsize{11}{12}\selectfont
\setlength{\parskip}{0pt}
\setlength{\abovedisplayskip}{3pt}
\setlength{\belowdisplayskip}{3pt}

\noindent\rule{\linewidth}{0.8pt}\par\nobreak

\begin{algorithmic}[1]

\Require Outer-training partition
$\mathcal D_{ec}^{\mathrm{t_{out}}}$,
frozen molecular logits
$\boldsymbol\eta^{\mathrm{chem}}_{ec}$,
and ridge grid $\Lambda$.

\Ensure Selected ridge penalty
$\widehat\lambda_{ec}$, or failure if no candidate is valid
across all three inner folds.

\Function{SelectLambda}
{$\mathcal D_{ec}^{\mathrm{t_{out}}},\Lambda$}

    \State Split $\mathcal D_{ec}^{\mathrm{t_{out}}}$ into three
    scaffold-disjoint inner folds.

    \For{each inner fold $\ell=1,2,3$}

        \State Define fold $\ell$ as
        $\mathcal D_{ec}^{\mathrm{val},\ell}$
        and the remaining folds as
        $\mathcal D_{ec}^{\mathrm{t_{in}},\ell}$.

        \State Fit
        $\widehat a_{ec}^{\mathrm{t_{in}},\ell}$
        on $\mathcal D_{ec}^{\mathrm{t_{in}},\ell}$,
        with the molecular logit coefficient fixed to one
        (Eq.~\ref{eq:calibration-intercept}).

        \State Estimate
        $\boldsymbol\mu_{ec}^{\mathrm{t_{in}},\ell}$ and
        $\mathbf d_{ec}^{\mathrm{t_{in}},\ell}$
        from $\mathcal D_{ec}^{\mathrm{t_{in}},\ell}$,
        and apply them unchanged to
        $\mathcal D_{ec}^{\mathrm{val},\ell}$
        (Eq.~\ref{eq:standardization}).

        \For{each $\lambda\in\Lambda$}

            \State With
            $\widehat a_{ec}^{\mathrm{t_{in}},\ell}$ fixed, fit
            $\widehat{\boldsymbol\beta}_{ec}^{\mathrm{t_{in}},\ell}
            (\lambda)$
            on $\mathcal D_{ec}^{\mathrm{t_{in}},\ell}$
            (Eq.~\ref{eq:residual-objective}).

            \State Evaluate the fitted model on
            $\mathcal D_{ec}^{\mathrm{val},\ell}$ and store
            $\mathcal L_{ec}^{\mathrm{val},\ell}(\lambda)$
            (Eq.~\ref{eq:inner-validation-loss});
            mark $\lambda$ invalid if fitting fails.

        \EndFor

    \EndFor

    \State Retain only penalties valid across all three inner folds,
    yielding $\Lambda_{\mathrm{valid}}$.

    \If{$\Lambda_{\mathrm{valid}}=\varnothing$}
        \State \Return Failure.
    \EndIf

    \State For each $\lambda\in\Lambda_{\mathrm{valid}}$, compute
    $\overline{\mathcal L}_{ec}^{\mathrm{val}}(\lambda)$
    across the three inner folds
    (Eq.~\ref{eq:mean-validation-loss}).

    \State Determine
    $\lambda_{ec}^{\min}$,
    the candidate minimizing mean validation loss
    (Eq.~\ref{eq:lambda-min}).

    \State Compute
    $\operatorname{SE}_{ec}(\lambda_{ec}^{\min})$
    across the three inner-validation losses
    (Eq.~\ref{eq:validation-se}).

    \State Select the largest valid penalty within one standard error
    of the minimum,
    $\widehat\lambda_{ec}$
    (Eq.~\ref{eq:lambda-selection}).

    \State \Return $\widehat\lambda_{ec}$.

\EndFunction

\end{algorithmic}

\par\nobreak\noindent\rule{\linewidth}{0.8pt}

\end{table}
\endgroup
\end{landscape}
\clearpage
% ============================================================
% Table S8
% ============================================================
\begin{landscape}
\begingroup

\begin{table}[p]
\centering

\captionsetup{
    font=small,
    width=\linewidth,
    skip=4pt
}

\caption{\textbf{Effect of molecular-intercept calibration.}
Changes are calibrated minus uncalibrated molecular predictions for
6-h cohorts. Values summarize 20 repetitions; Nadeau-Bengio (NB)
confidence intervals and tests use 100 paired outer-fold evaluations.
Negative $\Delta$LogLoss indicates improved calibration. AUPRC is
unchanged because intercept calibration preserves molecular ranking. \textbf{ALT TEXT:}
Table comparing calibrated and uncalibrated molecular predictions for four endpoints (NR-Aromatase, SR-ARE, SR-MMP, and SR-p53) across four 6-hour cell-line
cohorts. Log loss decreases after calibration in every cohort, while
AUPRC remains unchanged.}
\label{tab:s1-calibration-selected-endpoints}

\fontsize{8.5}{9.5}\selectfont
\setlength{\tabcolsep}{1.8pt}
\setlength{\extrarowheight}{0pt}
\renewcommand{\arraystretch}{1.05}

\resizebox{\linewidth}{!}{%
\begin{tabular}{@{}llrrrrrrrr@{}}

\toprule
\multirow[c]{2}{*}{Endpoint} &
\multirow[c]{2}{*}{Cell} &
\multicolumn{4}{c}{Log loss} &
\multicolumn{4}{c}{AUPRC} \\
\cmidrule(lr){3-6}\cmidrule(lr){7-10}
& &
$\Delta\,\pm$ SD &
95\% CI &
$P$ &
$q$ &
$\Delta\,\pm$ SD &
95\% CI &
$P$ &
$q$ \\
\midrule

\multirow[c]{4}{*}{Aromatase}
& HA1E
& \(-0.2353\,{\scriptstyle\pm\,0.0077}\)
& \([-0.2858,\,-0.1849]\)
& \(4.73\times10^{-15}\)
& \(1.42\times10^{-14}\)
& \(0.0000\,{\scriptstyle\pm\,0.0000}\)
& \([0.0000,\,0.0000]\)
& \textit{--}
& \textit{--} \\

& HepG2
& \(-0.2234\,{\scriptstyle\pm\,0.0100}\)
& \([-0.2796,\,-0.1672]\)
& \(4.30\times10^{-12}\)
& \(1.29\times10^{-11}\)
& \(0.0000\,{\scriptstyle\pm\,0.0000}\)
& \([0.0000,\,0.0000]\)
& \textit{--}
& \textit{--} \\

& HT29
& \(-0.2336\,{\scriptstyle\pm\,0.0075}\)
& \([-0.2885,\,-0.1787]\)
& \(2.70\times10^{-13}\)
& \(5.41\times10^{-13}\)
& \(0.0000\,{\scriptstyle\pm\,0.0000}\)
& \([0.0000,\,0.0000]\)
& \textit{--}
& \textit{--} \\

& MCF7
& \(-0.2372\,{\scriptstyle\pm\,0.0067}\)
& \([-0.2769,\,-0.1975]\)
& \(1.04\times10^{-20}\)
& \(3.11\times10^{-20}\)
& \(0.0000\,{\scriptstyle\pm\,0.0000}\)
& \([0.0000,\,0.0000]\)
& \textit{--}
& \textit{--} \\

\midrule

\multirow[c]{4}{*}{ARE}
& HA1E
& \(-0.0736\,{\scriptstyle\pm\,0.0078}\)
& \([-0.1122,\,-0.0350]\)
& \(2.62\times10^{-4}\)
& \(2.85\times10^{-4}\)
& \(0.0000\,{\scriptstyle\pm\,0.0000}\)
& \([0.0000,\,0.0000]\)
& \textit{--}
& \textit{--} \\

& HepG2
& \(-0.0648\,{\scriptstyle\pm\,0.0088}\)
& \([-0.1055,\,-0.0240]\)
& \(2.10\times10^{-3}\)
& \(2.30\times10^{-3}\)
& \(0.0000\,{\scriptstyle\pm\,0.0000}\)
& \([0.0000,\,0.0000]\)
& \textit{--}
& \textit{--} \\

& HT29
& \(-0.0829\,{\scriptstyle\pm\,0.0058}\)
& \([-0.1181,\,-0.0476]\)
& \(9.70\times10^{-6}\)
& \(1.16\times10^{-5}\)
& \(0.0000\,{\scriptstyle\pm\,0.0000}\)
& \([0.0000,\,0.0000]\)
& \textit{--}
& \textit{--} \\

& MCF7
& \(-0.0843\,{\scriptstyle\pm\,0.0094}\)
& \([-0.1238,\,-0.0449]\)
& \(4.89\times10^{-5}\)
& \(5.34\times10^{-5}\)
& \(0.0000\,{\scriptstyle\pm\,0.0000}\)
& \([0.0000,\,0.0000]\)
& \textit{--}
& \textit{--} \\

\midrule

\multirow[c]{4}{*}{MMP}
& HA1E
& \(-0.0820\,{\scriptstyle\pm\,0.0063}\)
& \([-0.1117,\,-0.0524]\)
& \(3.02\times10^{-7}\)
& \(4.03\times10^{-7}\)
& \(0.0000\,{\scriptstyle\pm\,0.0000}\)
& \([0.0000,\,0.0000]\)
& \textit{--}
& \textit{--} \\

& HepG2
& \(-0.0577\,{\scriptstyle\pm\,0.0067}\)
& \([-0.0928,\,-0.0225]\)
& \(1.60\times10^{-3}\)
& \(1.90\times10^{-3}\)
& \(0.0000\,{\scriptstyle\pm\,0.0000}\)
& \([0.0000,\,0.0000]\)
& \textit{--}
& \textit{--} \\

& HT29
& \(-0.0864\,{\scriptstyle\pm\,0.0074}\)
& \([-0.1239,\,-0.0489]\)
& \(1.43\times10^{-5}\)
& \(1.56\times10^{-5}\)
& \(0.0000\,{\scriptstyle\pm\,0.0000}\)
& \([0.0000,\,0.0000]\)
& \textit{--}
& \textit{--} \\

& MCF7
& \(-0.0940\,{\scriptstyle\pm\,0.0069}\)
& \([-0.1302,\,-0.0578]\)
& \(1.30\times10^{-6}\)
& \(1.56\times10^{-6}\)
& \(0.0000\,{\scriptstyle\pm\,0.0000}\)
& \([0.0000,\,0.0000]\)
& \textit{--}
& \textit{--} \\

\midrule

\multirow[c]{4}{*}{p53}
& HA1E
& \(-0.1603\,{\scriptstyle\pm\,0.0068}\)
& \([-0.1973,\,-0.1232]\)
& \(1.28\times10^{-13}\)
& \(3.07\times10^{-13}\)
& \(0.0000\,{\scriptstyle\pm\,0.0000}\)
& \([0.0000,\,0.0000]\)
& \textit{--}
& \textit{--} \\

& HepG2
& \(-0.1529\,{\scriptstyle\pm\,0.0079}\)
& \([-0.2026,\,-0.1033]\)
& \(1.95\times10^{-8}\)
& \(3.90\times10^{-8}\)
& \(0.0000\,{\scriptstyle\pm\,0.0000}\)
& \([0.0000,\,0.0000]\)
& \textit{--}
& \textit{--} \\

& HT29
& \(-0.1613\,{\scriptstyle\pm\,0.0059}\)
& \([-0.1963,\,-0.1262]\)
& \(8.73\times10^{-15}\)
& \(3.49\times10^{-14}\)
& \(0.0000\,{\scriptstyle\pm\,0.0000}\)
& \([0.0000,\,0.0000]\)
& \textit{--}
& \textit{--} \\

& MCF7
& \(-0.1748\,{\scriptstyle\pm\,0.0068}\)
& \([-0.2102,\,-0.1394]\)
& \(2.97\times10^{-16}\)
& \(5.95\times10^{-16}\)
& \(0.0000\,{\scriptstyle\pm\,0.0000}\)
& \([0.0000,\,0.0000]\)
& \textit{--}
& \textit{--} \\

\bottomrule

\end{tabular}%
}

\end{table}

\endgroup
\end{landscape}
\clearpage
% ============================================================
% Table S9
% ============================================================
\begin{landscape}
\begingroup
\begin{table}[p]
\centering
\caption{\textbf{Performance comparisons of molecular-only and transcriptome-completed models across toxicity endpoints and biological contexts.} 
Molecular-only model (Chem.) and transcriptome-complemented model (Chem. + Bio.) use same test
compounds per endpoint, cell line, and exposure time and are compared in a paired manner. Superscripts 
mark the strongest BH threshold met: \(^{\dagger}q<0.10\), \(^{*}q<0.05\),
\(^{**}q<0.01\), adjusted across twelve endpoints within cell, exposure,
cohort, and variant. Bold marks positive gains at \(q<0.10\);
header \((\uparrow)\) marks endpoints with at least one such displayed gain.
\(^{\ddagger}\): fewer than 20 active compounds. Abbreviations: BH, Benjamini-Hochberg. \textbf{ALT TEXT:}
Table comparing molecular-only and transcriptome-complemented AUPRC for
twelve toxicity endpoints across HA1E, HepG2, HT29, and MCF7 cells at
6 and 24 hours, with significance markers for positive performance gains.}
\label{tab:main-auprc}
\fontsize{9}{11}\selectfont
\setlength{\tabcolsep}{1.5pt}
\renewcommand{\arraystretch}{1.18}
% The vertical divider separates context identifiers from endpoint values.
\begin{tabular*}{\linewidth}{@{\extracolsep{\fill}}ccc|*{12}{c}@{}} % chktex 44
\toprule
& \multirow[c]{2}{*}{Time} & \multirow[c]{2}{*}{Cell} &
\multicolumn{7}{c}{Nuclear receptor (NR)} &
\multicolumn{5}{c}{Stress response (SR)} \\
\cmidrule(lr){4-10}\cmidrule(lr){11-15}
& & & AR & AR-LBD & AhR & \textbf{Aromatase}\,(\(\uparrow\)) & ER & ER-LBD & PPAR-\(\gamma\) & \textbf{ARE}\,(\(\uparrow\)) & ATAD5 & HSE & \textbf{MMP}\,(\(\uparrow\)) & \textbf{p53}\,(\(\uparrow\)) \\
\midrule
\multirow[c]{7}{*}[-2.5pt]{\rotatebox[origin=c]{90}{\fontsize{11}{13}\selectfont\bfseries Chem.}} & \multirow[c]{4}{*}{6~h} & HA1E & 0.461 & 0.554 & 0.526 & 0.384 & 0.430 & 0.385 & 0.235 & 0.474 & 0.295 & 0.234 & 0.588 & 0.361 \\
 &  & HepG2 & 0.537 & 0.554 & 0.477 & 0.346 & 0.477 & 0.404 & 0.300\(^{\ddagger}\) & 0.452 & 0.372 & 0.269 & 0.616 & 0.344 \\
 &  & HT29 & 0.457 & 0.561 & 0.479 & 0.372 & 0.419 & 0.397 & 0.193 & 0.451 & 0.307 & 0.254 & 0.563 & 0.338 \\
 &  & MCF7 & 0.459 & 0.581 & 0.473 & 0.409 & 0.402 & 0.379 & 0.222 & 0.452 & 0.254 & 0.248 & 0.527 & 0.295 \\
\cmidrule(lr){2-15}
 & \multirow[c]{3}{*}{24~h} & HA1E & 0.432 & 0.503 & 0.527 & 0.397 & 0.441 & 0.411 & 0.241 & 0.469 & 0.311 & 0.231 & 0.571 & 0.350 \\
 &  & HT29 & 0.385 & 0.510\(^{\ddagger}\) & 0.521 & 0.371 & 0.409 & 0.375 & 0.278\(^{\ddagger}\) & 0.498 & 0.308 & 0.297 & 0.547 & 0.374 \\
 &  & MCF7 & 0.460 & 0.583 & 0.468 & 0.411 & 0.407 & 0.397 & 0.212 & 0.449 & 0.257 & 0.239 & 0.529 & 0.294 \\
\midrule
\multirow[c]{7}{*}[-2.5pt]{\rotatebox[origin=c]{90}{\fontsize{11}{13}\selectfont\bfseries Chem. + Bio.}} & \multirow[c]{4}{*}{6~h} & HA1E & 0.461 & 0.554 & 0.526 & 0.384 & 0.430 & 0.384 & 0.235 & 0.479 & 0.293 & 0.248 & \textbf{0.646}\(^{\dagger}\) & 0.404 \\
 &  & HepG2 & 0.537 & 0.554 & 0.482 & 0.356 & 0.477 & 0.404 & 0.299\(^{\ddagger}\) & \textbf{0.553}\(^{\dagger}\) & 0.368 & 0.252 & \textbf{0.686}\(^{\dagger}\) & 0.395 \\
 &  & HT29 & 0.457 & 0.561 & 0.479 & 0.378 & 0.419 & 0.398 & 0.192 & 0.452 & 0.316 & 0.251 & \textbf{0.643}\(^{*}\) & 0.427 \\
 &  & MCF7 & 0.459 & 0.585 & 0.472 & 0.431 & 0.402 & 0.418 & 0.222 & 0.459 & 0.260 & 0.245 & \textbf{0.635}\(^{**}\) & 0.397 \\
\cmidrule(lr){2-15}
 & \multirow[c]{3}{*}{24~h} & HA1E & 0.432 & 0.503 & 0.530 & 0.482 & 0.441 & 0.400 & 0.241 & 0.482 & 0.337 & 0.238 & \textbf{0.681}\(^{\dagger}\) & 0.381 \\
 &  & HT29 & 0.385 & 0.510\(^{\ddagger}\) & 0.521 & 0.371 & 0.409 & 0.369 & 0.277\(^{\ddagger}\) & 0.492 & 0.308 & 0.299 & 0.553 & 0.399 \\
 &  & MCF7 & 0.460 & 0.586 & 0.468 & \textbf{0.538}\(^{\dagger}\) & 0.407 & 0.429 & 0.212 & 0.463 & 0.319 & 0.254 & \textbf{0.659}\(^{**}\) & \textbf{0.508}\(^{**}\) \\
\bottomrule
\end{tabular*}
\end{table}
\endgroup
\end{landscape}

\paragraph{Text S1: Quality control and processing of perturbation transcriptomic profiles.}

Quality control was performed before replicate aggregation. Perturbation
units represented by fewer than two replicates were excluded, as were treated
units with a plate-wise replicate Pearson correlation below 0.30.
Dimethyl sulfoxide (DMSO) vehicle-control profiles were additionally excluded
when mitochondrial-, ribosomal-, or hemoglobin-associated expression exceeded
three standard deviations from the corresponding distribution.

Profiles sharing the same compound, cell line, and exposure time were
aggregated gene-wise by the median. Perturbation responses were then referenced
to plate-matched DMSO controls by subtracting the arithmetic mean of the
corresponding control profiles. The resulting 978-dimensional response
$\mathbf x_c\in\mathbb R^{978}$ for each compound in context $c$ constituted
the context-specific transcriptomic dataset $\mathcal D^{\mathrm{bio}}_c$.

\clearpage
\end{document}
